\documentclass[11pt,oneside]{book}
\usepackage{notes-preamble}

\allowdisplaybreaks
\setcounter{secnumdepth}{2}
\setcounter{tocdepth}{2}
\emergencystretch=3em

% Automatically numbered environments, with one shared counter per section.
\theoremstyle{mystyle}
\newtheorem{notedefinition}{Definition}[section]
\newtheorem{noteproposition}[notedefinition]{Proposition}
\newtheorem{notetheorem}[notedefinition]{Theorem}
\newtheorem{notelemma}[notedefinition]{Lemma}
\newtheorem{notecorollary}[notedefinition]{Corollary}
\newtheorem{noteapplication}[notedefinition]{Application}
\newtheorem{noteexercise}[notedefinition]{Exercise}
\newtheorem{notealgorithm}[notedefinition]{Algorithm}
\newtheorem*{noteremark}{Remark}
\newtheorem{noteexample}[notedefinition]{Example}

\tcolorboxenvironment{notedefinition}{boxrule=0pt,boxsep=0pt,colback=red!10,left=8pt,right=8pt,enhanced jigsaw,borderline west={2pt}{0pt}{red},sharp corners,before skip=10pt,after skip=10pt,breakable}
\tcolorboxenvironment{noteproposition}{boxrule=0pt,boxsep=0pt,colback=Orange!10,left=8pt,right=8pt,enhanced jigsaw,borderline west={2pt}{0pt}{Orange},sharp corners,before skip=10pt,after skip=10pt,breakable}
\tcolorboxenvironment{notetheorem}{boxrule=0pt,boxsep=0pt,colback=blue!10,left=8pt,right=8pt,enhanced jigsaw,borderline west={2pt}{0pt}{blue},sharp corners,before skip=10pt,after skip=10pt,breakable}
\tcolorboxenvironment{notelemma}{boxrule=0pt,boxsep=0pt,colback=Cyan!10,left=8pt,right=8pt,enhanced jigsaw,borderline west={2pt}{0pt}{Cyan},sharp corners,before skip=10pt,after skip=10pt,breakable}
\tcolorboxenvironment{notecorollary}{boxrule=0pt,boxsep=0pt,colback=violet!10,left=8pt,right=8pt,enhanced jigsaw,borderline west={2pt}{0pt}{violet},sharp corners,before skip=10pt,after skip=10pt,breakable}
\tcolorboxenvironment{noteapplication}{boxrule=0pt,boxsep=0pt,colback=Orange!8,left=8pt,right=8pt,enhanced jigsaw,borderline west={2pt}{0pt}{Orange},sharp corners,before skip=10pt,after skip=10pt,breakable}
\tcolorboxenvironment{noteexercise}{boxrule=0pt,boxsep=0pt,colback=gray!8,left=8pt,right=8pt,enhanced jigsaw,borderline west={2pt}{0pt}{gray},sharp corners,before skip=10pt,after skip=10pt,breakable}
\tcolorboxenvironment{notealgorithm}{boxrule=0pt,boxsep=0pt,colback=yellow!12,left=8pt,right=8pt,enhanced jigsaw,borderline west={2pt}{0pt}{Goldenrod},sharp corners,before skip=10pt,after skip=10pt,breakable}
\tcolorboxenvironment{noteremark}{boxrule=0pt,boxsep=0pt,blanker,borderline west={2pt}{0pt}{Green},left=8pt,right=8pt,before skip=10pt,after skip=10pt,breakable}
\tcolorboxenvironment{noteexample}{boxrule=0pt,boxsep=0pt,blanker,borderline west={2pt}{0pt}{Black},left=8pt,right=8pt,before skip=10pt,after skip=10pt,breakable}

\newcommand{\comp}[1]{c\!\left(#1\right)}
\newcommand{\odd}{\operatorname{odd}}
\newcommand{\rad}{\operatorname{rad}}
\newcommand{\diam}{\operatorname{diam}}
\newcommand{\girth}{\operatorname{girth}}
\newcommand{\val}{\operatorname{val}}
\newcommand{\deficiency}{\operatorname{def}}
\newcommand{\cl}{\operatorname{cl}}
\newcommand{\dist}{\operatorname{dist}}
\newcommand{\out}{\operatorname{out}}
\newcommand{\inflow}{\operatorname{in}}
\newcommand{\supp}{\operatorname{supp}}
\newcommand{\Nzero}{\mathbb{N}_0}

\title{\Huge \textbf{Graph Theory}}
\author{Wei Chen}
\date{16 July 2026}

\begin{document}
\maketitle
	\allowdisplaybreaks
	\pagenumbering{roman}
	\setcounter{page}{1}
	\begin{center}
		\Large\textbf{Preface}
    \end{center}~\
    This is a lecture note of the course \text{Graph Theory} taught in 2021 Spring by Prof. Mario Huang at the Chinese University of Hong Kong, Shenzhen. Most of the content is based on
the presentation of the instructor and written down in my understanding, so there are unavoidable mistakes and typos in the note. The note is later revised and updated using ChatGPT. Topics of the note include: Eulerian and Hamiltonian graphs, connectivity, trees and paths, network flows, matchings, planar graphs, and graph colouring. 
	\newpage
	\tableofcontents
  \pagenumbering{arabic}
\chapter{Definitions and Eulerian Graphs}

\section{Basic terminology}

\begin{notedefinition}
A finite simple undirected graph is a pair $G=(V,E)$, where $V$ is a finite set and
\[
  E\subseteq \binom{V}{2}.
\]
The elements of $V$ are \emph{vertices} and the elements of $E$ are \emph{edges}.  We write $V(G)$ and $E(G)$ when the graph must be named explicitly.  An edge with ends $v,w$ is denoted $vw$ or $wv$.  If $vw\in E(G)$, then $v$ and $w$ are \emph{adjacent}; each is a \emph{neighbour} of the other.  The degree $\deg_G(v)$ is the number of edges incident with $v$.
\end{notedefinition}

\begin{notetheorem}[Handshaking lemma]
For every finite graph $G=(V,E)$,
\[
  \sum_{v\in V}\deg_G(v)=2|E|.
\]
\end{notetheorem}
\begin{proof}
Count vertex--edge incidences in two ways.  Every edge contributes one incidence at each of its two ends.
\end{proof}

\begin{notecorollary}
Every finite graph has an even number of odd-degree vertices.
\end{notecorollary}
\begin{proof}
The sum of all degrees is even, so the number of odd summands is even.
\end{proof}

\begin{notedefinition}
A \emph{multigraph} is a triple $(V,E,\varphi)$, where $V$ and $E$ are sets and
\[
  \varphi:E\longrightarrow \bigl\{\{u,v\}:u,v\in V\bigr\}
\]
assigns two ends to every edge.  If $\varphi(e)=\{v,v\}$, then $e$ is a \emph{loop}.  Distinct edges having the same pair of ends are \emph{parallel}.
\end{notedefinition}
\begin{notedefinition}
Let $G=(V,E)$.
\begin{itemize}
  \item A \emph{walk} is a sequence $v_0,e_1,v_1,\ldots,e_m,v_m$ with $e_i=v_{i-1}v_i$.  Its length is $m$.
  \item A walk with no repeated edge is a \emph{trail}.
  \item A walk with no repeated vertex is a \emph{path}.
  \item A walk is \emph{closed} if $v_0=v_m$.
  \item A closed trail is a \emph{circuit}.
  \item A closed walk of length at least three in which $v_0,\ldots,v_{m-1}$ are distinct is a \emph{cycle}.
\end{itemize}
\end{notedefinition}

\begin{notedefinition}
The minimum and maximum degrees of $G$ are
\[
  \delta(G)=\min_{v\in V(G)}\deg_G(v),
  \qquad
  \Delta(G)=\max_{v\in V(G)}\deg_G(v).
\]
\end{notedefinition}

\begin{notetheorem}
If $\delta(G)\ge 2$, then $G$ contains a cycle.
\end{notetheorem}
\begin{proof}
Choose $v_0\in V(G)$.  Having chosen $v_0,\ldots,v_i$, choose a neighbour $v_{i+1}$ of $v_i$ different from $v_{i-1}$; this is possible because $\deg(v_i)\ge2$.  Since $G$ is finite, some vertex eventually repeats.  At the first repetition, say $v_k=v_j$ with $j<k$, the vertices $v_j,v_{j+1},\ldots,v_{k-1}$ are distinct.  Hence
\[
  v_jv_{j+1}\cdots v_{k-1}v_j
\]
is a cycle.
\end{proof}
\begin{notedefinition}
A graph $H$ is a \emph{subgraph} of $G$ if $V(H)\subseteq V(G)$ and $E(H)\subseteq E(G)$.  For $X\subseteq V(G)$, the \emph{induced subgraph} $G[X]$ has vertex set $X$ and all edges of $G$ with both ends in $X$.

A graph is \emph{connected} if every two vertices are joined by a path.  A \emph{component} is a maximal connected subgraph.
\end{notedefinition}

\begin{notedefinition}
An \emph{isomorphism} from $G$ to $H$ is a bijection $f:V(G)\to V(H)$ such that
\[
  uv\in E(G)\quad\Longleftrightarrow\quad f(u)f(v)\in E(H).
\]
If such a map exists, then $G$ and $H$ are \emph{isomorphic}, written $G\cong H$.
\end{notedefinition}

\begin{notedefinition}
For vertices $v,w$ in a connected graph $G$, the distance $d_G(v,w)$ is the length of a shortest $v$--$w$ path.  If no such path exists, set $d_G(v,w)=\infty$.
\end{notedefinition}

\begin{notedefinition}
A set $X\subseteq V(G)$ is \emph{independent} if no two vertices of $X$ are adjacent.  A graph is \emph{bipartite} if its vertex set can be partitioned into two independent sets $A$ and $B$; we write $G=(A,B)$.  The complete bipartite graph with parts of sizes $s$ and $t$ is denoted $K_{s,t}$.
\end{notedefinition}
\begin{notetheorem}
A graph is bipartite if and only if it contains no odd cycle.
\end{notetheorem}
\begin{proof}
If $G=(A,B)$ is bipartite, every walk alternates between $A$ and $B$, so every cycle has even length.

Conversely, it is enough to treat one connected component at a time.  Fix a root $r$ and set
\[
  A=\{v:d(r,v)\text{ is even}\},
  \qquad
  B=\{v:d(r,v)\text{ is odd}\}.
\]
Suppose an edge $uv$ has both ends in the same set.  Let $P_u$ and $P_v$ be shortest paths from $r$ to $u$ and $v$, and let $x$ be their last common vertex.  The two tails from $x$ to $u$ and $v$ are internally disjoint and have lengths of the same parity.  Together with $uv$ they form an odd cycle, a contradiction.  Thus $A$ and $B$ are independent, and $G$ is bipartite.
\end{proof}
\section{Eulerian graphs and bridges}

\begin{notedefinition}
A trail containing every edge of $G$ is an \emph{Eulerian trail}.  A closed Eulerian trail is an \emph{Eulerian circuit}.  A graph is \emph{Eulerian} if it has an Eulerian circuit and \emph{semi-Eulerian} if it has an Eulerian trail but no Eulerian circuit.
\end{notedefinition}

\begin{notedefinition}
For $F\subseteq E(G)$, the graph $G-F$ is obtained by deleting the edges in $F$.  For $S\subseteq V(G)$, the graph $G-S$ is obtained by deleting the vertices in $S$ and every edge incident with them.
\end{notedefinition}

\begin{notetheorem}[Euler's theorem]
A connected graph $G$ is Eulerian if and only if every vertex has even degree.
\end{notetheorem}
\begin{proof}
In an Eulerian circuit, every visit to a vertex uses one entering and one leaving edge, so every degree is even.

Conversely, assume every degree is even.  Start at any vertex and follow unused edges until no continuation is possible.  The trail cannot stop at a vertex different from its starting vertex, because at every other visited vertex the used incident edges occur in entering--leaving pairs.  Hence the trail is closed.  If it does not yet contain every edge, connectedness gives a vertex of the circuit incident with an unused edge.  Starting there produces another closed trail, which can be spliced into the first.  Repeating this finite process yields an Eulerian circuit.
\end{proof}

\begin{notecorollary}
A connected graph with at least two vertices is semi-Eulerian if and only if it has exactly two odd-degree vertices.  Every Eulerian trail starts at one odd vertex and ends at the other.
\end{notecorollary}
\begin{proof}
In any trail, all internal vertices use incident edges in pairs; only the two ends can have odd degree.  Conversely, if $u$ and $v$ are the only odd vertices, add one extra edge $uv$ (allowing a parallel edge if necessary).  The resulting multigraph is Eulerian.  Deleting the added edge from an Eulerian circuit leaves an Eulerian trail from $u$ to $v$.
\end{proof}

\begin{notedefinition}
A set $F\subseteq E(G)$ is a \emph{disconnecting set} if $G-F$ has more components than $G$.  A minimal nonempty disconnecting set is an \emph{edge cut}.  An edge $e$ is a \emph{bridge} if $\{e\}$ is an edge cut.
\end{notedefinition}

\begin{notelemma}
If there is a walk between two vertices, then there is a path between them.
\end{notelemma}
\begin{proof}
Repeatedly delete the segment between two consecutive occurrences of a repeated vertex.  The process terminates with a walk having no repeated vertices.
\end{proof}

\begin{notedefinition}
The complement $\overline G$ of a simple graph $G$ has vertex set $V(G)$ and
\[
  uv\in E(\overline G)
  \quad\Longleftrightarrow\quad
  u\ne v\text{ and }uv\notin E(G).
\]
A clique in $G$ is an independent set in $\overline G$; equivalently, it is a set of pairwise adjacent vertices.
\end{notedefinition}
\begin{notetheorem}
Let $G$ be connected and let $e=uv\in E(G)$.
\begin{enumerate}
  \item The edge $e$ is a bridge if and only if it lies on no cycle.
  \item If $e$ is a bridge, then $G-e$ has exactly two components, one containing $u$ and one containing $v$.
\end{enumerate}
\end{notetheorem}
\begin{proof}
If $e$ lies on a cycle, the rest of that cycle still joins $u$ to $v$ after $e$ is removed; replacing every use of $e$ by this alternate path shows that $G-e$ remains connected.  Thus $e$ is not a bridge.

Conversely, if $e$ is not a bridge, then $G-e$ contains a $u$--$v$ path.  Adding $e$ to that path produces a cycle containing $e$.

Now assume $e$ is a bridge.  The vertices $u$ and $v$ lie in different components of $G-e$.  For any vertex $x$, choose an $x$--$u$ path in $G$.  If it avoids $e$, then $x$ lies in the $u$-component of $G-e$; if it uses $e$, the portion before or after $e$ places $x$ in the $v$-component.  Hence there are exactly two components.
\end{proof}
\begin{notedefinition}
A proper vertex colouring is a map $c:V(G)\to S$ such that adjacent vertices receive different colours.  The chromatic number $\chi(G)$ is the minimum size of a colour set admitting a proper colouring.
\end{notedefinition}

\begin{notetheorem}
If every vertex of a connected graph $G$ has even degree, then $G$ has no bridge.
\end{notetheorem}
\begin{proof}
Suppose $e=uv$ were a bridge, and let $A$ be the component of $G-e$ containing $u$.  In the graph $G[A]$, the vertex $u$ has odd degree and every other vertex has even degree.  This contradicts the handshaking lemma.
\end{proof}

\begin{notealgorithm}[Fleury's algorithm]
Let $G$ be connected and all its degrees even.
\begin{enumerate}
  \item Choose a starting vertex $v_0$.
  \item At the current vertex, traverse an unused incident edge that is not a bridge of the current unused-edge graph whenever such an edge exists.  Use a bridge only when it is the sole available continuation.
  \item Delete each traversed edge from the unused-edge graph and continue until no edge remains.
\end{enumerate}
The resulting closed trail is an Eulerian circuit.
\end{notealgorithm}
\begin{proof}
At every stage, all vertices of the unused-edge graph have even degree except possibly the starting vertex and the current vertex, which have odd degree when they are distinct.  Consequently the current vertex cannot become stranded away from the start.  Moreover, the rule against crossing a bridge preserves the component containing all unused edges that can still be reached from the current position; a bridge is used only when no unused edge would be abandoned on the current side.  Hence the process returns to $v_0$ only after all edges have been used.  Since the graph is finite, the algorithm terminates and produces an Eulerian circuit.
\end{proof}
\section{Further definitions and elementary structure}

\begin{notedefinition}
\begin{itemize}
  \item $G$ is \emph{self-complementary} if $G\cong\overline G$.
  \item The \emph{girth} $g(G)$ is the length of a shortest cycle; an acyclic graph has infinite girth.
  \item An automorphism is an isomorphism $G\to G$.  A graph is \emph{vertex-transitive} if its automorphism group acts transitively on $V(G)$.
  \item The union $G_1\cup\cdots\cup G_k$ has vertex set $\bigcup_iV(G_i)$ and edge set $\bigcup_iE(G_i)$.
  \item If $G$ and $H$ have disjoint vertex sets, their \emph{join} $G+H$ is obtained from $G\cup H$ by adding every edge between $V(G)$ and $V(H)$.
  \item A graph is $H$-free if it has no induced subgraph isomorphic to $H$.
  \item A nonnegative integer sequence is \emph{graphic} if it is the degree sequence of a simple graph.
\end{itemize}
\end{notedefinition}
\begin{notetheorem}
The complete graph $K_n$ is the union of $k$ bipartite spanning subgraphs if and only if $n\le 2^k$.
\end{notetheorem}
\begin{proof}
Assume first that $n\le2^k$.  Label the vertices by distinct binary strings of length $k$.  For each coordinate $i$, let $H_i$ contain every edge whose endpoints differ in coordinate $i$.  The graph $H_i$ is bipartite, with parts given by bit $0$ and bit $1$.  Any two distinct strings differ in some coordinate, so every edge of $K_n$ lies in at least one $H_i$.

Conversely, suppose $K_n=H_1\cup\cdots\cup H_k$ with each $H_i$ bipartite.  Add isolated vertices to each $H_i$ if necessary and choose a bipartition $(A_i,B_i)$.  Give a vertex $v$ the binary string whose $i$th bit records whether $v\in A_i$ or $v\in B_i$.  Distinct vertices must receive distinct strings: their edge belongs to some $H_i$, where they lie on opposite sides.  Hence $n\le2^k$.
\end{proof}

\begin{notetheorem}
If every vertex of a simple graph $G$ has degree at least $k$, then $G$ contains a path of length at least $k$.  If $k\ge2$, then $G$ also contains a cycle of length at least $k+1$.
\end{notetheorem}
\begin{proof}
Let $P=v_0v_1\cdots v_m$ be a longest path.  Every neighbour of $v_0$ lies on $P$, so $m\ge\deg(v_0)\ge k$.  Let $v_i$ be the neighbour of $v_0$ with largest index.  At least $k$ distinct neighbours occur among $v_1,\ldots,v_i$, so $i\ge k$.  Therefore
\[
  v_0v_1\cdots v_iv_0
\]
is a cycle of length $i+1\ge k+1$.
\end{proof}
\begin{notetheorem}
Let $G$ be a connected nontrivial graph with exactly $2k$ odd-degree vertices.  The minimum number of trails whose edge sets partition $E(G)$ is $\max\{k,1\}$.
\end{notetheorem}
\begin{proof}
Every trail has either zero or two odd endpoints.  Thus $k$ trails are necessary to account for $2k$ odd vertices; a nontrivial graph needs at least one trail.

For the upper bound when $k\ge1$, pair the odd vertices and add one auxiliary edge for each pair.  All degrees become even, so the resulting connected multigraph has an Eulerian circuit.  Deleting the $k$ auxiliary edges breaks that circuit into at most $k$ trails covering the original edges.  When $k=0$, one Eulerian circuit suffices.
\end{proof}

\begin{notetheorem}[Havel--Hakimi]
Let $d=(d_1,\ldots,d_n)$ be a nonincreasing sequence of nonnegative integers with $d_1\le n-1$.  Let $d'$ be obtained by deleting $d_1$, subtracting $1$ from the next $d_1$ entries, and reordering.  Then $d$ is graphic if and only if $d'$ is graphic.
\end{notetheorem}
\begin{proof}
If $d'$ is realized by a graph on $v_2,\ldots,v_n$, add a new vertex $v_1$ adjacent to the $d_1$ vertices whose degrees were reduced.

Conversely, take a realization of $d$ and let $v_1$ have degree $d_1$.  Among all such realizations, choose one in which $v_1$ is adjacent to as many as possible of $v_2,\ldots,v_{d_1+1}$.  If $v_1v_i$ is missing for some $i\le d_1+1$, then $v_1v_j$ is present for some $j>d_1+1$.  Since $\deg(v_i)\ge\deg(v_j)$, there is a vertex $x$ adjacent to $v_i$ but not to $v_j$.  Replacing the edges $v_1v_j$ and $v_ix$ by $v_1v_i$ and $v_jx$ is a valid $2$-switch and increases the chosen number, a contradiction.  Thus $v_1$ is adjacent to $v_2,\ldots,v_{d_1+1}$, and deleting $v_1$ realizes $d'$.
\end{proof}
\begin{notedefinition}[2-switch]
If $xy$ and $uv$ are edges of a simple graph while $xu$ and $yv$ are nonedges, replacing $xy,uv$ by $xu,yv$ is a \emph{$2$-switch}.  It preserves every vertex degree.
\end{notedefinition}

\begin{notetheorem}
Two simple graphs on the same labelled vertex set have the same degree at every vertex if and only if one can be transformed into the other by a sequence of $2$-switches.
\end{notetheorem}
\begin{proof}
A $2$-switch preserves all degrees.  Conversely, colour the edges of $G-H$ red and the edges of $H-G$ blue.  At each vertex the red and blue degrees are equal, so the symmetric difference decomposes into closed even alternating trails.  On a shortest nontrivial alternating trail, an appropriate pair of red edges can be switched to two blue edges; shortestness guarantees that the inserted pair is absent from the current graph.  This strictly decreases $|E(G)\triangle E(H)|$.  Induction completes the transformation.
\end{proof}
\begin{notedefinition}
The eccentricity of $v$ is $\max_x d(v,x)$.  A vertex is \emph{central} if its eccentricity is minimum.  The radius and diameter are
\[
  \rad(G)=\min_v\max_x d(v,x),
  \qquad
  \diam(G)=\max_{u,v}d(u,v).
\]
For every connected graph,
\[
  \rad(G)\le\diam(G)\le2\rad(G).
\]
\end{notedefinition}

\begin{noteproposition}
Every graph containing a cycle satisfies
\[
  g(G)\le2\diam(G)+1.
\]
\end{noteproposition}
\begin{proof}
Let $C=v_0v_1\cdots v_{g-1}v_0$ be a shortest cycle.  The shorter arc of $C$ between any two of its vertices is a shortest path in $G$, for otherwise a shortcut would create a shorter cycle.  Hence
\[
  \left\lfloor\frac g2\right\rfloor\le\diam(G),
\]
which is equivalent to the stated bound.
\end{proof}

\begin{noteproposition}
If $G$ has radius at most $k$ and maximum degree at most $d\ge3$, then
\[
 |V(G)|\le 1+d\sum_{i=0}^{k-1}(d-1)^i
 =\frac{d(d-1)^k-2}{d-2}
 <\frac d{d-2}(d-1)^k.
\]
\end{noteproposition}
\begin{proof}
Choose a central vertex $v_0$ and let $S_i=\{v:d(v_0,v)=i\}$.  Then $|S_0|=1$, $|S_1|\le d$, and, for $i\ge1$, every vertex of $S_i$ has at most $d-1$ neighbours in $S_{i+1}$.  Thus $|S_i|\le d(d-1)^{i-1}$ for $i\ge1$.  Summing through distance $k$ gives the bound.
\end{proof}
\section{Selected exercises}

\begin{noteexercise}
Let $G$ be connected and let $F$ be an edge cut.  Prove that:
\begin{enumerate}
  \item $G-F$ has exactly two components;
  \item for every $e=uv\in F$, the vertices $u$ and $v$ lie in different components of $G-F$.
\end{enumerate}
\end{noteexercise}
\begin{proof}
Fix $e\in F$.  By minimality of the edge cut, $G-(F\setminus\{e\})$ is connected, and $e$ is a bridge of that graph.  Theorem~1.8 therefore shows that deleting $e$ leaves exactly two components.  But
\[
  \bigl(G-(F\setminus\{e\})\bigr)-e=G-F,
\]
so both claims follow.
\end{proof}
\begin{noteexercise}
For a connected graph $G$, prove that $G$ is bipartite if and only if, for every fixed vertex $r$, no two adjacent vertices have the same distance from $r$.
\end{noteexercise}
\begin{proof}
If $G$ is bipartite, the parity of $d(r,v)$ records the side of the bipartition containing $v$, so adjacent vertices have opposite parities and hence different distances.

Conversely, suppose $G$ is not bipartite and choose a shortest odd cycle
\[
  C=v_0v_1\cdots v_{2k}v_0.
\]
The two vertices $v_k$ and $v_{k+1}$ are adjacent.  We claim that both have distance $k$ from $v_0$.  The cycle gives paths of length $k$ to each.  If, say, a shorter $v_0$--$v_k$ path existed, combining it with one of the two arcs of $C$ would yield an odd cycle shorter than $C$.  Thus
\[
  d(v_0,v_k)=d(v_0,v_{k+1})=k,
\]
contradicting the stated distance property.
\end{proof}
\begin{noteexercise}
Show that every graph $G$ has a map $\alpha:V(G)\to\{0,1\}$ such that, at every vertex $v$, at least half of the edges incident with $v$ join vertices receiving different values.
\end{noteexercise}
\begin{proof}
Choose a partition $V(G)=A\sqcup B$ maximizing the number of crossing edges.  If a vertex $v\in A$ had fewer than half of its incident edges crossing to $B$, moving $v$ to $B$ would increase the cut size, a contradiction.  The same argument applies to vertices of $B$.  Set $\alpha=0$ on $A$ and $\alpha=1$ on $B$.
\end{proof}
\begin{noteremark}[References recorded in the source]
D.~B.~West, \emph{Introduction to Graph Theory}, Chapter~1; R.~J.~Wilson, \emph{Introduction to Graph Theory}, Chapter~1 and Section~2.2; \emph{Graph Theory with Applications}, Chapter~1 and Section~4.2; R.~Diestel, \emph{Graph Theory}, Chapter~1.
\end{noteremark}

\chapter{Hamiltonian Graphs and Digraphs}

\section{Hamiltonian cycles and closure}

\begin{notedefinition}
A \emph{Hamiltonian path} is a path containing every vertex of $G$.  A \emph{Hamiltonian cycle} is a cycle containing every vertex.  A graph is \emph{Hamiltonian} if it has a Hamiltonian cycle and \emph{semi-Hamiltonian} if it has a Hamiltonian path but no Hamiltonian cycle.
\end{notedefinition}

\begin{notetheorem}
If $G$ is Hamiltonian, then for every nonempty $S\subseteq V(G)$,
\[
  \comp{G-S}\le |S|.
\]
\end{notetheorem}
\begin{proof}
Let $C$ be a Hamiltonian cycle.  Deleting $k=|S|$ vertices from a cycle leaves at most $k$ path components.  Since $C-S$ is a spanning subgraph of $G-S$, adding the remaining edges of $G-S$ cannot increase the number of components.
\end{proof}

\begin{notedefinition}
For $F\subseteq\binom{V(G)}2$, the graph $G+F$ has vertex set $V(G)$ and edge set $E(G)\cup F$.  We write $G+uv$ when $F=\{uv\}$.
\end{notedefinition}

\begin{notedefinition}
Let $G$ have $n$ vertices.  Starting with $G_0=G$, repeatedly add an edge $uv$ between nonadjacent vertices whenever
\[
  \deg_{G_i}(u)+\deg_{G_i}(v)\ge n.
\]
Stop when no such pair remains.  The terminal graph is the \emph{Hamiltonian closure} of $G$, denoted $\cl(G)$.
\end{notedefinition}
\begin{notetheorem}
The Hamiltonian closure of a graph is unique; it does not depend on the order in which eligible edges are added.
\end{notetheorem}
\begin{proof}
Call a supergraph $H$ of $G$ \emph{closed} if every nonedge $uv$ of $H$ satisfies $\deg_H(u)+\deg_H(v)<n$.  A terminal graph produced by the procedure is closed.

Let $K$ be any closed supergraph of $G$.  We show inductively that every edge added during any closure sequence belongs to $K$.  Suppose the current graph $H$ is contained in $K$ and $uv$ is eligible in $H$.  If $uv\notin E(K)$, then
\[
  \deg_K(u)+\deg_K(v)\ge \deg_H(u)+\deg_H(v)\ge n,
\]
contradicting that $K$ is closed.  Thus $uv\in E(K)$.  Hence every terminal graph is contained in every closed supergraph of $G$, and two terminal graphs contain one another.
\end{proof}

\begin{notelemma}[Ore's edge-addition lemma]
Let $G$ have $n$ vertices, and let $u,v$ be nonadjacent with
\[
  \deg_G(u)+\deg_G(v)\ge n.
\]
Then $G$ is Hamiltonian if and only if $G+uv$ is Hamiltonian.
\end{notelemma}
\begin{proof}
Only the reverse implication needs proof.  Let $C$ be a Hamiltonian cycle in $G+uv$.  If $C$ does not use $uv$, then $C\subseteq G$.  Otherwise deleting $uv$ from $C$ gives a Hamiltonian path
\[
  P=v_1v_2\cdots v_n,
  \qquad v_1=u,
  \quad v_n=v.
\]
Define
\[
  A=\{i:uv_{i+1}\in E(G)\},
  \qquad
  B=\{i:v_iv\in E(G)\},
\]
where $1\le i\le n-1$.  We have $|A|=\deg(u)$ and $|B|=\deg(v)$.  Since $|A|+|B|\ge n>n-1$, the sets intersect.  Choose $i\in A\cap B$.  Then
\[
  v_1v_2\cdots v_i v_n v_{n-1}\cdots v_{i+1}v_1
\]
is a Hamiltonian cycle in $G$.
\end{proof}
\begin{notetheorem}[Bondy--Chv\'atal]
A graph $G$ is Hamiltonian if and only if $\cl(G)$ is Hamiltonian.
\end{notetheorem}
\begin{proof}
Apply Lemma~2.3 after each edge addition in a closure sequence.
\end{proof}

\begin{notecorollary}
If $|V(G)|\ge3$ and $\cl(G)$ is complete, then $G$ is Hamiltonian.
\end{notecorollary}

\begin{notetheorem}[Ore's theorem]
Let $G$ have $n\ge3$ vertices.  If
\[
  \deg(u)+\deg(v)\ge n
\]
for every nonadjacent pair $u,v$, then $G$ is Hamiltonian.
\end{notetheorem}
\begin{proof}
Every nonedge is eligible, so $\cl(G)=K_n$.  Apply Corollary~2.5.
\end{proof}

\begin{notetheorem}[Dirac's theorem]
If $G$ has $n\ge3$ vertices and $\delta(G)\ge n/2$, then $G$ is Hamiltonian.
\end{notetheorem}
\begin{proof}
Every nonadjacent pair has degree sum at least $n$, so Ore's theorem applies.
\end{proof}

\begin{notetheorem}[Chv\'atal's degree-sequence condition]
Let $G$ have $n\ge3$ vertices and nondecreasing degree sequence
\[
  d_1\le d_2\le\cdots\le d_n.
\]
Suppose that for every integer $i<n/2$,
\[
  d_i\le i\quad\Longrightarrow\quad d_{n-i}\ge n-i.
\]
Then $G$ is Hamiltonian.
\end{notetheorem}
\begin{proof}
It is enough to show that $H=\cl(G)$ is complete.  The Chv\'atal condition is preserved when edges are added, so $H$ has the same property.  Suppose $H$ is not complete.  Choose a nonedge $uv$ maximizing $\deg_H(u)+\deg_H(v)$, with $\deg_H(u)\le\deg_H(v)$, and put $i=\deg_H(u)$.  Since $H$ is closed,
\[
  2i\le\deg_H(u)+\deg_H(v)\le n-1,
\]
so $i<n/2$.

Every nonneighbour of $v$ has degree at most $i$, by maximality of the chosen degree sum.  Since $v$ has at least $i$ nonneighbours, $d_i(H)\le i$.  Likewise every nonneighbour of $u$ has degree at most $\deg_H(v)\le n-1-i$; together with $u$, this gives at least $n-i$ vertices of degree at most $n-1-i$.  Hence
\[
  d_{n-i}(H)\le n-i-1<n-i,
\]
contradicting the degree-sequence hypothesis.  Thus $H$ is complete, and Bondy--Chv\'atal gives the result.
\end{proof}
\section{Tournaments and kernels}

\begin{notedefinition}
A digraph is a pair $D=(V,A)$, where $A$ is a set of ordered pairs of distinct vertices, called \emph{arcs}.  We write $uv$ for the arc $(u,v)$.  A directed walk, path, or cycle follows the directions of its arcs.  The underlying graph is obtained by forgetting all directions.  The digraph is \emph{connected} if its underlying graph is connected and \emph{strongly connected} if every ordered pair of vertices is joined by a directed path.  A \emph{tournament} is an orientation of a complete graph: for each pair $u\ne v$, exactly one of $uv$ and $vu$ is an arc.
\end{notedefinition}

\begin{notetheorem}
Every tournament has a directed Hamiltonian path.
\end{notetheorem}
\begin{proof}
Use induction on $n=|V(D)|$.  Let
\[
  P=v_1v_2\cdots v_{n-1}
\]
be a directed Hamiltonian path in the tournament obtained by deleting $x$.  If $xv_1$ is an arc, place $x$ first; if $v_{n-1}x$ is an arc, place $x$ last.  Otherwise there is a first index $i$ such that $xv_{i+1}$ is an arc.  By minimality, $v_ix$ is an arc, and
\[
  v_1\cdots v_i x v_{i+1}\cdots v_{n-1}
\]
is a directed Hamiltonian path.
\end{proof}
\begin{notetheorem}[Camion's theorem for tournaments]
Every strongly connected tournament with at least three vertices has a directed Hamiltonian cycle.
\end{notetheorem}
\begin{proof}
Let $C$ be a longest directed cycle.  Suppose that $R=V(D)\setminus V(C)$ is nonempty.  If some $x\in R$ has an in-neighbour and an out-neighbour on $C$, then two consecutive vertices $v_i,v_{i+1}$ of $C$ satisfy
\[
  v_i x,\; xv_{i+1}\in A(D),
\]
and $x$ can be inserted into $C$, a contradiction.  Therefore every $x\in R$ either dominates all of $C$ or is dominated by all of $C$.

Let
\[
  A=\{x\in R:C\to x\},
  \qquad
  B=\{x\in R:x\to C\}.
\]
Strong connectivity implies that both sets are nonempty and that some arc goes from $A$ to $B$: take a directed path from a vertex of $A$ to $C$ and look at the first transition from $A$ to $B$.  Choose $a\in A$, $b\in B$ with $ab\in A(D)$.  Replacing any arc $v_iv_{i+1}$ of $C$ by
\[
  v_i\to a\to b\to v_{i+1}
\]
creates a longer directed cycle, again a contradiction.  Hence $C$ is spanning.
\end{proof}
\begin{notedefinition}
A set $K\subseteq V(D)$ is a \emph{kernel} if it is independent and every vertex in $V(D)\setminus K$ has an outgoing arc to a vertex of $K$.
\end{notedefinition}

\begin{notetheorem}[Richardson]
Every finite digraph with no directed odd cycle has a kernel.
\end{notetheorem}
\begin{proof}
We use induction on $|V(D)|$.  First suppose $D$ is strongly connected.  Fix a root $r$.  Because every directed closed walk decomposes into directed cycles, every directed closed walk has even length.  Consequently, any two directed $r$--$v$ paths have the same parity: append a directed $v$--$r$ path to both and compare the resulting closed walks.  Partition the vertices according to that parity.  Every arc goes between the two parts, so the underlying graph is bipartite.  Since strong connectivity gives every vertex positive outdegree, either part is a kernel.

Now suppose $D$ is not strongly connected.  Let $C$ be a sink strongly connected component and let $K_C$ be a kernel of $D[C]$.  Let
\[
  X=\{x\notin C:\text{ some arc }xk\text{ has }k\in K_C\}.
\]
By induction, the induced digraph $D-(C\cup X)$ has a kernel $K'$.  Then $K_C\cup K'$ is independent: no arc leaves the sink component $C$, and a vertex of $K'$ has no arc to $K_C$ by the definition of $X$.  Vertices in $C\setminus K_C$ are absorbed by $K_C$, vertices in $X$ are absorbed by $K_C$, and all remaining vertices are absorbed by $K'$.  Thus $K_C\cup K'$ is a kernel of $D$.
\end{proof}
\begin{notedefinition}
The \emph{split} of a digraph $D$ is the bipartite graph with parts
\[
  V^+=\{v^+:v\in V(D)\},
  \qquad
  V^-=\{v^-:v\in V(D)\},
\]
and with edge $u^+v^-$ for every arc $uv$ of $D$.
\end{notedefinition}
\section{Eulerian digraphs and applications}

\begin{notetheorem}
A finite digraph $D$ is Eulerian if and only if
\[
  d^+(v)=d^-(v)\qquad\text{for every }v\in V(D)
\]
and its underlying graph has at most one nontrivial component.
\end{notetheorem}
\begin{proof}
The conditions are necessary because an Eulerian circuit enters and leaves each vertex equally often and contains every nonisolated vertex.

Conversely, begin at any nonisolated vertex and follow unused outgoing arcs.  At any vertex other than the start, the number of used incoming arcs exceeds the number of used outgoing arcs by one, so equality of indegree and outdegree guarantees an unused outgoing arc.  Thus the trail closes.  If unused arcs remain, connectedness of the nontrivial underlying component gives a vertex on the current circuit incident with an unused arc.  Build another closed directed trail there and splice it into the circuit.  Repetition yields an Eulerian circuit.
\end{proof}

\begin{noteapplication}[de Bruijn cycles]
For $n\ge1$, form the digraph whose vertices are the binary strings of length $n-1$ and whose arcs are the binary strings $a_1\cdots a_n$, directed from $a_1\cdots a_{n-1}$ to $a_2\cdots a_n$.  Every vertex has indegree and outdegree $2$, and the underlying graph has one nontrivial component.  An Eulerian circuit therefore lists all $2^n$ binary strings of length $n$ exactly once as consecutive cyclic blocks.  Such a cyclic word is a binary de Bruijn cycle of order $n$.
\end{noteapplication}
\begin{notedefinition}[King]
A vertex $v$ of a digraph is a \emph{king} if every vertex is reachable from $v$ by a directed path of length at most two.
\end{notedefinition}

\begin{noteproposition}[Landau]
Every tournament has a king.
\end{noteproposition}
\begin{proof}
Choose a vertex $v$ of maximum outdegree.  Let $u$ be a vertex not dominated by $v$, so $u\to v$.  If no out-neighbour $w$ of $v$ satisfies $w\to u$, then $u$ dominates $v$ and every out-neighbour of $v$.  Hence
\[
  d^+(u)\ge d^+(v)+1,
\]
a contradiction.  Thus $v\to w\to u$ for some $w$, and $v$ is a king.
\end{proof}
\begin{notetheorem}[double-tree bound]
Let $G$ be a complete weighted graph whose weights are nonnegative and satisfy the triangle inequality
\[
  w(xz)\le w(xy)+w(yz).
\]
If $T$ is a minimum spanning tree and $C^*$ is a minimum-weight Hamiltonian cycle, then there is a Hamiltonian cycle $C$ such that
\[
  w(C)\le2w(T)\le2w(C^*).
\]
\end{notetheorem}
\begin{proof}
Duplicate every edge of $T$.  The resulting connected multigraph has all degrees even, so it has an Eulerian circuit $W$ of weight $2w(T)$.  Traverse $W$ and shortcut every visit to a vertex already seen, retaining only the first occurrence of each vertex before returning to the start.  Completeness supplies every shortcut edge, and the triangle inequality ensures that shortcutting never increases total weight.  The result is a Hamiltonian cycle $C$ with $w(C)\le2w(T)$.

Deleting any edge from $C^*$ gives a spanning tree, so $w(T)\le w(C^*)$.
\end{proof}
\begin{noteremark}[References recorded in the source]
D.~B.~West, \emph{Introduction to Graph Theory}, Section~7.2; R.~J.~Wilson, \emph{Introduction to Graph Theory}, Section~2.3; \emph{Graph Theory with Applications}, Section~4.2; R.~Diestel, \emph{Graph Theory}, Chapter~10.
\end{noteremark}

\chapter{Connectivity}

\section{Vertex and edge connectivity}

\begin{notedefinition}
A set $X\subseteq V(G)$ is a \emph{separating set} if $G-X$ has more components than $G$.  A separating set consisting of one vertex is a \emph{cut-vertex}.
\end{notedefinition}

\begin{notedefinition}
The vertex connectivity $\kappa(G)$ is the least integer $k$ for which there is a set $X$ of $k$ vertices such that $G-X$ is disconnected or has at most one vertex.  The graph is $k$-connected if $\kappa(G)\ge k$.
\end{notedefinition}

\begin{noteremark}
\[
  \kappa(G)=0
  \quad\Longleftrightarrow\quad
  G\text{ is disconnected or }|V(G)|=1,
\]
and $\kappa(K_n)=n-1$.  Every connected graph with at least two vertices is $1$-connected.  A graph is $k$-connected if and only if it is $m$-connected for every $m\le k$.
\end{noteremark}

\begin{notedefinition}
The edge connectivity $\lambda(G)$ is the least integer $k$ for which there is a set $F$ of $k$ edges such that $G-F$ is disconnected or has only one vertex.  The graph is $k$-edge-connected if $\lambda(G)\ge k$.
\end{notedefinition}

\begin{notetheorem}
Let $G$ be connected and let $F$ be an edge cut.  Then $G-F$ has exactly two components.  Moreover, for every $e=uv\in F$, the vertices $u$ and $v$ lie in different components of $G-F$.
\end{notetheorem}
\begin{proof}
Fix $e\in F$.  Minimality of $F$ implies that $G-(F\setminus\{e\})$ is connected, and $e$ is a bridge of that graph.  The bridge theorem therefore implies that deleting $e$ produces exactly two components with $u$ and $v$ on opposite sides.  The resulting graph is precisely $G-F$.
\end{proof}
\begin{notetheorem}[Whitney's inequalities]
For every nontrivial simple graph $G$,
\[
  \kappa(G)\le\lambda(G)\le\delta(G).
\]
\end{notetheorem}
\begin{proof}
Deleting all edges incident with a minimum-degree vertex disconnects that vertex from the rest, so $\lambda(G)\le\delta(G)$.  Deleting all neighbours of such a vertex also disconnects it (or leaves a graph with at most one vertex), so $\kappa(G)\le\delta(G)$.

It remains to prove $\kappa(G)\le\lambda(G)$.  The claim is immediate when $G$ is disconnected.  Let $F$ be a minimum edge cut of a connected graph, and let $A,B$ be the two components of $G-F$.  If one component is a singleton $\{v\}$, then
\[
  \lambda(G)=\deg(v)\ge\delta(G),
\]
so $\lambda(G)=\delta(G)$ and $\kappa(G)\le\lambda(G)$.

Assume $|A|,|B|\ge2$, and fix $e=uv\in F$ with $u\in A$ and $v\in B$.  For each edge of $F\setminus\{e\}$ choose one endpoint different from $u$ and $v$, and let $S$ be the set of chosen endpoints.  Then $|S|\le|F|-1$.  In $G-S$, the only possible edge between $A-S$ and $B-S$ is $uv$.  If $G-S$ is disconnected, then $S$ is already a vertex cut.  Otherwise $uv$ is a bridge of $G-S$; deleting $u$ or $v$ (choosing the end that does not remove an entire side) disconnects the graph.  Thus $G$ has a vertex cut of size at most $|S|+1\le|F|$.
\end{proof}
\section{Menger's theorem}

\begin{notedefinition}
Let $A,B\subseteq V(G)$.
\begin{itemize}
  \item An $(A,B)$-path has one end in $A$, one end in $B$, and no internal vertex in $A\cup B$.
  \item An $(A,B)$-separator is a set $S\subseteq V(G)$ such that $G-S$ contains no $(A,B)$-path.  A path of length $0$ at a vertex of $A\cap B$ is allowed, so every separator contains $A\cap B$.
  \item An $(A,B)$-connector is a set of pairwise vertex-disjoint $(A,B)$-paths.
\end{itemize}
\end{notedefinition}

\begin{notetheorem}[Menger's theorem, vertex form]
For every finite graph $G$ and all $A,B\subseteq V(G)$, the maximum size of an $(A,B)$-connector equals the minimum size of an $(A,B)$-separator.
\end{notetheorem}
\begin{proof}
Every separator must meet every path in a connector, so the minimum separator size is at least the maximum connector size.

For the reverse inequality, use the vertex-splitting reduction to integral max flow.  Replace each vertex $x$ by an arc $x^-x^+$ of capacity $1$, replace every edge $xy$ by the two arcs $x^+y^-$ and $y^+x^-$ of sufficiently large capacity, add a source joined to $a^-$ for $a\in A$, and join $b^+$ to a sink for $b\in B$.  Give the source and sink arcs sufficiently large capacity, except that vertices in $A\cap B$ are handled by their compulsory length-$0$ paths.  Integral flows correspond exactly to families of vertex-disjoint $(A,B)$-paths, while finite cuts correspond exactly to vertex separators.  The max-flow min-cut theorem, proved in Chapter~5, gives equality.
\end{proof}
\begin{notedefinition}
For distinct vertices $u,v$, two $u$--$v$ paths are \emph{internally disjoint} if their vertex sets intersect only in $\{u,v\}$.
\end{notedefinition}

\begin{notecorollary}
If $u$ and $v$ are distinct nonadjacent vertices, then the minimum size of a $u$--$v$ separator equals the maximum number of internally disjoint $u$--$v$ paths.
\end{notecorollary}

\begin{notecorollary}
A graph with at least two vertices is $k$-connected if and only if every two distinct vertices are joined by $k$ internally disjoint paths.
\end{notecorollary}

\begin{notetheorem}
A graph with at least three vertices is $2$-connected if and only if every two vertices lie on a common cycle.
\end{notetheorem}
\begin{proof}
Two internally disjoint $u$--$v$ paths have a union containing a cycle through $u$ and $v$.  Conversely, the two arcs of a cycle through $u,v$ are internally disjoint $u$--$v$ paths.
\end{proof}

\begin{notetheorem}[edge form of Menger]
For distinct vertices $u,v$, the minimum size of an edge set whose deletion destroys all $u$--$v$ paths equals the maximum number of pairwise edge-disjoint $u$--$v$ paths.
\end{notetheorem}

\begin{notetheorem}
A graph with at least two vertices is $k$-edge-connected if and only if every two distinct vertices are joined by $k$ edge-disjoint paths.
\end{notetheorem}
\section{Cubic graphs}

\begin{notedefinition}
For $S,T\subseteq V(G)$, write $[S,T]$ for the set of edges with one end in $S$ and the other in $T$.  An edge cut has the form $[S,\overline S]$ for a nonempty proper subset $S\subset V(G)$.
\end{notedefinition}

\begin{notetheorem}
If $G$ is $3$-regular, then
\[
  \kappa(G)=\lambda(G).
\]
\end{notetheorem}
\begin{proof}
Whitney's inequalities give $\kappa(G)\le\lambda(G)\le3$.

If $\lambda(G)=0$, then both connectivities are $0$.  If $\lambda(G)=1$, a bridge exists.  Removing an endpoint on a nontrivial side of the bridge disconnects $G$, so $\kappa(G)=1$.

Suppose $\lambda(G)=2$.  If $G$ had a cut-vertex $v$, then the three edges incident with $v$ would be distributed among at least two components of $G-v$; some component would receive only one of them, and that edge would be a bridge.  Hence $\kappa(G)\ne1$, so $\kappa(G)=2$.

Finally suppose $\lambda(G)=3$ and, for contradiction, $\kappa(G)=2$.  Let $\{u,v\}$ be a vertex cut.  Every component of $G-\{u,v\}$ sends at least three edges to $\{u,v\}$, because $\lambda(G)=3$.  Since $u$ and $v$ have only six incident edge slots in total, there are exactly two components, $uv\notin E(G)$, and each component sends exactly three edges to $\{u,v\}$.  Neither component can send all three edges to the same one of $u,v$, since that would make the other cut vertex a cut-vertex and force a bridge.  Thus each component has a $1$--$2$ distribution.  The distributions are opposite on the two components.  Deleting the two edges that are unique at their respective cut vertices separates the graph into two parts, producing a $2$-edge cut, contrary to $\lambda(G)=3$.  Therefore $\kappa(G)=3$.
\end{proof}
\begin{noteexercise}[Fan lemma]
For a simple graph $G$, prove that the following are equivalent:
\begin{enumerate}
  \item $G$ is $k$-connected;
  \item $|V(G)|>k$, and for every $v\in V(G)$ and every $S\subseteq V(G)\setminus\{v\}$ with $|S|\ge k$, there are $k$ paths from $v$ to $S$ that intersect only at $v$.
\end{enumerate}
\end{noteexercise}
\begin{proof}
Assume $G$ is $k$-connected.  Add a new vertex $x$ adjacent to every vertex of $S$.  Any $v$--$x$ separator in the new graph has size at least $k$: deleting fewer than $k$ vertices leaves $G$ connected and leaves at least one vertex of $S$.  Menger's theorem gives $k$ internally disjoint $v$--$x$ paths.  Deleting $x$ yields the required $v$--$S$ fan, with distinct endpoints in $S$.

Conversely, suppose the fan condition holds but $X$ is a separating set with $|X|<k$.  Choose $v$ in one component of $G-X$ and $w$ in another.  Choose a $k$-set $S$ containing $X\cup\{w\}$ and not containing $v$.  A $k$-fan from $v$ to $S$ has one path ending at every vertex of $S$.  The path ending at $w$ must pass through some $x\in X$; it then meets the distinct path ending at $x$, contradicting that the fan paths intersect only at $v$.
\end{proof}
\begin{noteremark}[References recorded in the source]
D.~B.~West, \emph{Introduction to Graph Theory}, Chapter~4; R.~J.~Wilson, \emph{Introduction to Graph Theory}, Sections~2.1 and~6.2; \emph{Graph Theory with Applications}, Chapter~3; R.~Diestel, \emph{Graph Theory}, Chapter~3.
\end{noteremark}

\chapter{Trees}

\section{Trees and spanning trees}

\begin{notedefinition}
A graph is \emph{acyclic} if it contains no cycle.  An acyclic graph is a \emph{forest}; a connected forest is a \emph{tree}.  A vertex of degree one in a forest is a \emph{leaf}.
\end{notedefinition}

\begin{notetheorem}
Let $T$ be a graph with $n$ vertices.  The following statements are equivalent.
\begin{enumerate}
  \item $T$ is a tree.
  \item $T$ is acyclic and has $n-1$ edges.
  \item $T$ is connected and has $n-1$ edges.
  \item $T$ is connected and every edge is a bridge.
  \item Every two vertices are joined by a unique path.
  \item $T$ is acyclic, and $T+uv$ contains exactly one cycle for every nonedge $uv$.
\end{enumerate}
\end{notetheorem}
\begin{proof}
A longest path in a nontrivial tree has leaf endpoints.  Removing a leaf and its incident edge leaves a smaller tree, so induction gives $|E(T)|=n-1$.  Thus (a) implies (b).

If an acyclic graph has components of orders $n_1,\ldots,n_k$, each component is a tree and has $n_i-1$ edges; hence it has $n-k$ edges in total.  Therefore (b) implies $k=1$, giving (c).

Every connected $n$-vertex graph has at least $n-1$ edges: start with one vertex and expose the remaining vertices along paths from the initial set, adding at least one new edge for every new vertex.  If a connected graph with exactly $n-1$ edges contained a cycle, deleting one cycle edge would leave a connected graph with only $n-2$ edges.  Hence (c) implies (a).

In a connected graph, an edge is a bridge exactly when it lies on no cycle, so (a) and (d) are equivalent.  A connected graph has two distinct $u$--$v$ paths exactly when their union contains a cycle, giving the equivalence with (e).  Finally, in an acyclic graph the unique $u$--$v$ path, together with a new edge $uv$, creates exactly one cycle; this proves the equivalence with (f).
\end{proof}

\begin{notecorollary}
Every tree with at least two vertices has at least two leaves.
\end{notecorollary}
\begin{proof}
Let $P=v_0v_1\cdots v_m$ be a longest path.  Any neighbour of $v_0$ outside $P$ would extend the path, and any second neighbour on $P$ would create a cycle.  Thus $\deg(v_0)=1$.  Similarly $\deg(v_m)=1$.
\end{proof}
\begin{notedefinition}
A subgraph $H$ of $G$ is \emph{spanning} if $V(H)=V(G)$.  A spanning subgraph that is a tree is a \emph{spanning tree}.
\end{notedefinition}

\begin{notetheorem}
Every connected graph contains a spanning tree.
\end{notetheorem}
\begin{proof}
Choose a tree $T\subseteq G$ with as many vertices as possible.  If $T$ is not spanning, connectedness gives an edge $uv$ with $u\in V(T)$ and $v\notin V(T)$.  Then $T+uv$ is a larger tree, a contradiction.  Hence $T$ spans $G$.
\end{proof}

\begin{noteremark}[Equivalent deletion proof]
Repeatedly delete one edge from a cycle.  Such a deletion does not disconnect the graph.  The finite process ends with a connected acyclic spanning subgraph, hence a spanning tree.
\end{noteremark}
\begin{notetheorem}[spanning-tree exchange]
Let $T$ and $T'$ be distinct spanning trees of $G$.  For every $e\in E(T)\setminus E(T')$, there is an edge $e'\in E(T')\setminus E(T)$ such that
\[
  T-e+e'
\]
is a spanning tree.
\end{notetheorem}
\begin{proof}
The graph $T-e$ has exactly two components, with vertex sets $A$ and $B$.  The unique path in $T'$ between the ends of $e$ must cross from $A$ to $B$, so it contains an edge $e'\in[A,B]$.  This edge is not in $T$, since $e$ is the only edge of $T$ crossing that cut.  Adding $e'$ reconnects the two components of $T-e$ without creating a cycle.
\end{proof}
\section{Minimum spanning trees}

\begin{notedefinition}
Let $G$ have an edge-weight function $w:E(G)\to\mathbb R$.  The weight of a subgraph $H$ is
\[
  w(H)=\sum_{e\in E(H)}w(e).
\]
A \emph{minimum spanning tree} (MST) is a spanning tree of minimum total weight.
\end{notedefinition}

\begin{notealgorithm}[Kruskal]
For a connected weighted graph $G$:
\begin{enumerate}
  \item Start with the spanning forest $T=(V(G),\varnothing)$.
  \item Repeatedly choose a minimum-weight edge joining two different components of $T$, and add it to $T$.
  \item Stop after $n-1$ edges have been selected.
\end{enumerate}
The output is an MST.
\end{notealgorithm}
\begin{proof}
Maintain the invariant that the selected forest $T_i$ is contained in some MST $M_i$.  The invariant is trivial initially.  Let $e$ be the next selected edge.  If $e\in M_i$, do nothing.  Otherwise $M_i+e$ contains a cycle.  That cycle contains an edge $f\notin T_i$ crossing between two components of $T_i$.  Kruskal's choice gives $w(e)\le w(f)$, so
\[
  M_{i+1}=M_i-f+e
\]
is again an MST and contains $T_i+e$.  At termination the selected graph itself is a spanning tree contained in an MST, so it is an MST.
\end{proof}
\begin{notealgorithm}[Prim]
For a connected weighted graph $G$:
\begin{enumerate}
  \item Choose a root $v_0$, set $W=\{v_0\}$, and let $T$ have no edges.
  \item While $W\ne V(G)$, choose a minimum-weight edge $uv$ with $u\in W$ and $v\notin W$; add $uv$ to $T$ and add $v$ to $W$.
\end{enumerate}
The output is an MST.
\end{notealgorithm}
\begin{proof}
Again maintain an MST $M$ containing the current tree $T$.  If the selected edge $e$ is not in $M$, then $M+e$ contains a cycle.  This cycle contains another edge $f$ crossing the cut $(W,V\setminus W)$.  Since Prim selected a lightest crossing edge, $w(e)\le w(f)$, and $M-f+e$ is an MST containing the enlarged $T$.  At the end, $T$ is spanning and therefore is an MST.
\end{proof}
\section{Shortest paths}

\begin{notedefinition}
Let $D=(V,A)$ be a digraph with arc weights $w:A\to\mathbb R$.  The weight of a directed path is the sum of its arc weights.  The distance $d(u,v)$ is the minimum weight of a directed $u$--$v$ path, or $\infty$ if no such path exists.  If a negative directed cycle is reachable from $u$ and can reach $v$, then the infimum is $-\infty$; throughout the finite-distance results below, such a cycle is excluded.
\end{notedefinition}

\begin{notealgorithm}[Dijkstra]
Assume all arc weights are nonnegative and choose a source $s$.
\begin{enumerate}
  \item Set $W=\varnothing$, $\ell(s)=0$, and $\ell(v)=\infty$ for $v\ne s$.
  \item While $W\ne V$, choose $x\in V\setminus W$ with minimum $\ell(x)$, add $x$ to $W$, and for every arc $xy$ set
  \[
    \ell(y)\leftarrow\min\{\ell(y),\ell(x)+w(xy)\}.
  \]
\end{enumerate}
When the algorithm ends, $\ell(v)=d(s,v)$ for every $v$.
\end{notealgorithm}
\begin{proof}
We prove that whenever a vertex $x$ is added to $W$, its label is final: $\ell(x)=d(s,x)$.  Suppose all earlier vertices satisfy the claim.  A finite label is always the weight of an $s$--$x$ walk, so $\ell(x)\ge d(s,x)$.

Let $P$ be a shortest $s$--$x$ path, and let $y$ be its first vertex outside the old set $W$, with predecessor $z\in W$.  By induction, $\ell(z)=d(s,z)$, so relaxation of $zy$ gave
\[
  \ell(y)\le d(s,z)+w(zy)=d(s,y)\le d(s,x),
\]
where the last inequality uses nonnegative weights on the remaining part of $P$.  Since $x$ was chosen with minimum label outside $W$,
\[
  \ell(x)\le\ell(y)\le d(s,x).
\]
Thus equality holds.
\end{proof}

\begin{noteremark}[Running time]
A direct array implementation uses $O(|V|^2+|A|)$ operations.  With a binary heap the usual bound is $O((|V|+|A|)\log|V|)$.  Nonnegativity, not strict positivity, is the required hypothesis; negative arcs invalidate the greedy finalization step.
\end{noteremark}
\begin{notedefinition}
For sets $S_1,\ldots,S_k$, their disjoint union is
\[
  S_1\sqcup\cdots\sqcup S_k
  =\{(x,i):x\in S_i,\ 1\le i\le k\}.
\]
\end{notedefinition}

\begin{notetheorem}
The multiset of arcs of any directed $u$--$v$ walk can be written as the disjoint union of the arcs of a directed $u$--$v$ path and the arcs of zero or more directed cycles.
\end{notetheorem}
\begin{proof}
If the walk has no repeated vertex, it is already a path.  Otherwise choose two occurrences of the same vertex with no repeated vertex between them.  The intervening segment is a directed cycle.  Delete that segment and apply induction to the shorter $u$--$v$ walk.
\end{proof}
\begin{notetheorem}[triangle inequality]
If $D$ has no negative directed cycle, then for all vertices $x,y,z$,
\[
  d(x,y)\le d(x,z)+d(z,y),
\]
with the usual conventions for $\infty$.
\end{notetheorem}
\begin{proof}
Concatenate a shortest $x$--$z$ path with a shortest $z$--$y$ path.  The resulting walk decomposes into an $x$--$y$ path and directed cycles.  Every directed cycle has nonnegative weight, so the path has weight at most the concatenated walk.
\end{proof}
\begin{notealgorithm}[Bellman--Ford]
Let $D=(V,A)$ have no negative cycle reachable from a source $s$.
\begin{enumerate}
  \item Set $\ell(s)=0$ and $\ell(v)=\infty$ for $v\ne s$.
  \item Repeat $|V|-1$ times: for every arc $uv\in A$, perform the relaxation
  \[
    \ell(v)\leftarrow\min\{\ell(v),\ell(u)+w(uv)\}.
  \]
\end{enumerate}
Then $\ell(v)=d(s,v)$ for all $v$.
\end{notealgorithm}
\begin{proof}
Let $d_i(s,v)$ be the minimum weight of an $s$--$v$ path using at most $i$ arcs.  After the $i$th full pass,
\[
  \ell(v)\le d_i(s,v).
\]
Indeed, for a shortest path with at most $i$ arcs, its prefix has at most $i-1$ arcs, and when the final arc is processed its relaxation gives the desired bound.

Every finite label is the weight of an $s$--$v$ walk.  By Theorem~4.5 and the absence of negative cycles, such a walk has weight at least $d(s,v)$.  Hence $\ell(v)\ge d(s,v)$ at every stage.  A shortest path can be chosen simple and therefore has at most $|V|-1$ arcs.  Thus
\[
  d_{|V|-1}(s,v)=d(s,v),
\]
and the two inequalities force equality after the final pass.
\end{proof}
\begin{notetheorem}
Let $P$ be a shortest $s$--$u$ path, and let $v$ be the predecessor of $u$ on $P$.  Then the $s$--$v$ subpath of $P$ is a shortest $s$--$v$ path.
\end{notetheorem}
\begin{proof}
If a lighter $s$--$v$ path existed, concatenate it with the arc $vu$.  The resulting $s$--$u$ walk would be lighter than $P$.  Removing any directed cycles cannot increase its weight, so it would contain a lighter $s$--$u$ path, a contradiction.
\end{proof}

\begin{notecorollary}
Every subpath of a shortest path is shortest between its own endpoints.
\end{notecorollary}
\begin{notetheorem}[negative-cycle test]
After the usual $|V|-1$ Bellman--Ford passes, perform one additional full pass.  If a negative directed cycle is reachable from $s$, then at least one label decreases.
\end{notetheorem}
\begin{proof}
Let $C=v_0v_1\cdots v_{m-1}v_0$ be a reachable negative cycle.  After $|V|-1$ passes, every vertex of $C$ has a finite label.  Suppose no label decreased during the extra pass.  Then for every $i$ (indices modulo $m$),
\[
  \ell(v_{i+1})\le\ell(v_i)+w(v_iv_{i+1}).
\]
Summing gives
\[
  0\le w(C),
\]
contradicting that $C$ is negative.
\end{proof}
\section{Cut and cycle properties of minimum spanning trees}

\begin{notetheorem}
If all edge weights of a connected graph are distinct, then its MST is unique.
\end{notetheorem}
\begin{proof}
Suppose $T$ and $T'$ are distinct MSTs, and choose the minimum-weight edge $e\in E(T)\triangle E(T')$; without loss, $e\in T\setminus T'$.  Adding $e$ to $T'$ creates a cycle containing an edge $f\in T'\setminus T$.  By the choice of $e$ and distinctness of weights, $w(e)<w(f)$.  Then $T'-f+e$ is a lighter spanning tree, a contradiction.
\end{proof}

\begin{notetheorem}[cut property]
Let $S$ be a nonempty proper subset of $V(G)$.  If an edge $e\in[S,V\setminus S]$ has strictly smaller weight than every other edge crossing the cut, then every MST contains $e$.
\end{notetheorem}
\begin{proof}
Let $T$ be an MST not containing $e$.  Adding $e$ creates a cycle, which contains another cut edge $f$.  Then $w(e)<w(f)$ and $T-f+e$ is a lighter spanning tree, a contradiction.
\end{proof}

\begin{notetheorem}[cycle property]
If an edge $e$ is the unique maximum-weight edge on a cycle $C$, then no MST contains $e$.
\end{notetheorem}
\begin{proof}
If an MST $T$ contained $e$, then $T-e$ would have two components.  The remainder of $C$ contains an edge $f$ crossing between them, with $w(f)<w(e)$.  Thus $T-e+f$ would be a lighter spanning tree.
\end{proof}
\begin{notetheorem}
Let $T$ be an MST of a connected graph $G$, and suppose $v_0$ is a leaf of $T$.
\begin{enumerate}
  \item $G-v_0$ is connected.
  \item $T-v_0$ is an MST of $G-v_0$.
\end{enumerate}
\end{notetheorem}
\begin{proof}
For vertices $x,y\ne v_0$, the unique $x$--$y$ path in $T$ cannot use the leaf $v_0$ internally.  Hence it lies in $T-v_0$, proving that $G-v_0$ is connected and that $T-v_0$ is a spanning tree.

Let $e_0$ be the edge of $T$ incident with $v_0$.  If $G-v_0$ had a spanning tree $T'$ lighter than $T-v_0$, then $T'+e_0$ would be a spanning tree of $G$ with weight smaller than $T$, a contradiction.
\end{proof}
\begin{noteremark}[References recorded in the source]
\emph{Graph Theory with Applications}, Chapter~2; D.~B.~West, \emph{Introduction to Graph Theory}, Chapter~2; R.~J.~Wilson, \emph{Introduction to Graph Theory}, Chapter~3; additional source notes cite GTM~173, Chapters~1, 3, and~10.
\end{noteremark}

\chapter{Network Flows}

\section{Networks, flows, and cuts}

\begin{notedefinition}
A \emph{network} is a quadruple
\[
  N=(D,c,s,t),
\]
where $D=(V,A)$ is a finite directed graph, $c\colon A\to\R_{\ge0}$ is the
\emph{capacity function}, and $s,t\in V$ are distinct vertices called the
\emph{source} and the \emph{sink}.

For $X\subseteq V$ and a function $g\colon A\to\R$, put
\[
\begin{aligned}
  \delta^+(X)&:=\{xy\in A:x\in X,\ y\notin X\},\\
  \delta^-(X)&:=\{xy\in A:x\notin X,\ y\in X\},\\
  g_{\out}(X)&:=\sum_{e\in\delta^+(X)}g(e),
  &g_{\inflow}(X)&:=\sum_{e\in\delta^-(X)}g(e).
\end{aligned}
\]
For a single vertex $v$, we abbreviate $g_{\out}(\{v\})$ and
$g_{\inflow}(\{v\})$ to $g_{\out}(v)$ and $g_{\inflow}(v)$.
\end{notedefinition}

\begin{notedefinition}
A \emph{flow} in $N$ is a function $f\colon A\to\R_{\ge0}$ satisfying
\begin{enumerate}
  \item $0\le f(e)\le c(e)$ for every $e\in A$;
  \item $f_{\inflow}(v)=f_{\out}(v)$ for every $v\in V\setminus\{s,t\}$.
\end{enumerate}
The second condition is called \emph{flow conservation}.  The \emph{value} of
$f$ is
\[
  \val(f):=f_{\out}(s)-f_{\inflow}(s).
\]
The zero function is the \emph{zero flow}.
\end{notedefinition}

\begin{notedefinition}
A flow is \emph{maximum} if its value is at least the value of every other
flow in the same network.
\end{notedefinition}

\begin{notetheorem}
Let $f$ be a flow in $N=(D,c,s,t)$.  If $S\subseteq V$ satisfies
$s\in S$ and $t\notin S$, then
\[
  f_{\out}(S)-f_{\inflow}(S)=\val(f).
\]
In particular,
\[
  \val(f)=f_{\inflow}(t)-f_{\out}(t).
\]
\end{notetheorem}
\begin{proof}
Summing $f_{\out}(v)-f_{\inflow}(v)$ over $v\in S$ cancels every arc whose
two ends lie in $S$.  The uncancelled terms are exactly the arcs leaving or
entering $S$, so
\[
  f_{\out}(S)-f_{\inflow}(S)
  =\sum_{v\in S}\bigl(f_{\out}(v)-f_{\inflow}(v)\bigr).
\]
All summands vanish by flow conservation except the one at $s$, which is
$\val(f)$.  Taking $S=V\setminus\{t\}$ gives the second identity.
\end{proof}
\begin{notedefinition}
An $(s,t)$-\emph{cut} is a partition $V=S\dot\cup T$ with $s\in S$ and
$t\in T$.  We write the cut as $[S,T]$ and define its capacity by
\[
  c[S,T]:=c_{\out}(S)=\sum_{e\in\delta^+(S)}c(e).
\]
A cut of minimum capacity is a \emph{minimum cut}.
\end{notedefinition}

\begin{notetheorem}[weak duality]
For every flow $f$ and every $(s,t)$-cut $[S,T]$,
\[
  \val(f)\le c[S,T].
\]
\end{notetheorem}
\begin{proof}
By Theorem~5.1 and the capacity constraints,
\[
  \val(f)=f_{\out}(S)-f_{\inflow}(S)
  \le f_{\out}(S)
  \le c_{\out}(S)=c[S,T].
\]
\end{proof}
\section{Augmenting paths and residual networks}

\begin{notedefinition}
Let
\[
  P=(v_0,v_1,\dots,v_k),\qquad v_0=s,\quad v_k=t,
\]
be a path in the underlying undirected graph of $D$.  For each step
$v_{i-1}v_i$, call it
\begin{itemize}
  \item \emph{forward} if $v_{i-1}v_i\in A$, with residual capacity
        $c(v_{i-1}v_i)-f(v_{i-1}v_i)$;
  \item \emph{backward} if $v_iv_{i-1}\in A$, with residual capacity
        $f(v_iv_{i-1})$.
\end{itemize}
The residual capacity of $P$ is the minimum of these quantities.  If it is
positive, $P$ is an $f$-\emph{augmenting path}.  Augmenting by $\varepsilon$
means adding $\varepsilon$ to every forward arc of $P$ and subtracting
$\varepsilon$ from every backward arc.
\end{notedefinition}

\begin{noteremark}
An augmenting path records exactly where more flow can be pushed forward and
where existing flow can be cancelled and rerouted backward.
\end{noteremark}

\begin{notetheorem}
For a flow $f$ in a network $N$, the following are equivalent.
\begin{enumerate}
  \item $f$ is a maximum flow.
  \item There is no $f$-augmenting path from $s$ to $t$.
  \item There is an $(s,t)$-cut $[S,T]$ such that
        $c[S,T]=\val(f)$.
\end{enumerate}
\end{notetheorem}
\begin{proof}
If an augmenting path has residual capacity $\varepsilon>0$, augmenting along
it preserves the capacity constraints and flow conservation and increases the
flow value by $\varepsilon$.  Thus (1) implies (2).

Assume (2).  Form the residual digraph $D_f$: for each arc $xy\in A$, include
$xy$ when $c(xy)-f(xy)>0$, and include $yx$ when $f(xy)>0$.  Let $S$ be the
vertices reachable from $s$ in $D_f$, and let $T=V\setminus S$.  Since there
is no augmenting path, $t\in T$.  If $xy\in A$ has $x\in S$, $y\in T$, then
its forward residual capacity is zero, so $f(xy)=c(xy)$.  If $yx\in A$ has
$x\in S$, $y\in T$, then its backward residual capacity is zero, so
$f(yx)=0$.  Therefore
\[
\begin{aligned}
  \val(f)
  &=f_{\out}(S)-f_{\inflow}(S)\\
  &=\sum_{e\in\delta^+(S)}c(e)-0
   =c[S,T],
\end{aligned}
\]
and (3) holds.

Finally, (3) implies (1) by weak duality: every flow $g$ satisfies
$\val(g)\le c[S,T]=\val(f)$.
\end{proof}
\begin{notetheorem}[Max-Flow Min-Cut Theorem]
In every finite network, the maximum value of a flow equals the minimum
capacity of an $(s,t)$-cut.
\end{notetheorem}
\begin{proof}
A maximum flow exists because the feasible flows form a nonempty compact
polytope and the value is continuous.  Apply Theorem~5.3 to such a flow.
Weak duality gives the reverse inequality, so equality follows.
\end{proof}

\begin{notealgorithm}[Ford--Fulkerson]
Start with the zero flow.  While an augmenting path $P$ exists, augment by the
full residual capacity of $P$.  When no augmenting path remains, return the
current flow.
\end{notealgorithm}

\begin{notetheorem}
If every capacity is an integer, Ford--Fulkerson terminates after finitely many
augmentations and returns an integer-valued maximum flow.  In particular,
every network with integer capacities has an integer-valued maximum flow.
\end{notetheorem}
\begin{proof}
Starting from the zero flow, every residual capacity and every augmentation
amount is an integer.  Each augmentation increases the flow value by at least
$1$.  Weak duality bounds that value by the capacity of any fixed cut, so only
finitely many augmentations are possible.  At termination Theorem~5.3 gives a
maximum flow.
\end{proof}

\begin{noteremark}
For arbitrary real capacities, a poor choice of augmenting paths can lead to an
infinite sequence of augmentations.  The integer-capacity hypothesis is what
makes the elementary termination argument valid.
\end{noteremark}

\begin{notedefinition}
The \emph{residual digraph} $D_f$ of a flow $f$ has vertex set $V(D)$ and,
for every original arc $xy$, contains
\begin{itemize}
  \item a forward residual arc $xy$ of capacity $c(xy)-f(xy)$ when this is
        positive;
  \item a backward residual arc $yx$ of capacity $f(xy)$ when this is
        positive.
\end{itemize}
Parallel residual arcs may occur.  An augmenting path is precisely a directed
$s$--$t$ path in $D_f$.
\end{notedefinition}
\section{Applications}

\begin{noteapplication}[vertex-disjoint paths by vertex splitting]
Let $G$ be an undirected graph and let $u,v$ be distinct nonadjacent vertices.
Construct a network as follows.
\begin{enumerate}
  \item Keep $u$ as the source and $v$ as the sink.
  \item Replace every vertex $x\in V(G)\setminus\{u,v\}$ by two vertices
        $x^-$ and $x^+$ and the arc $x^-x^+$ of capacity $1$.
  \item For every edge $xy\in E(G)$, add arcs that represent traversal in
        both directions: from the output copy of one end to the input copy of
        the other.  Give these arcs a capacity larger than $|V(G)|$, or simply
        capacity $\infty$.
\end{enumerate}
Then the value of a maximum flow equals the maximum number of internally
vertex-disjoint $u$--$v$ paths.  By Max-Flow Min-Cut, it also equals the
minimum size of a $u$--$v$ vertex separator.
\end{noteapplication}
\begin{proof}
Every integral unit of flow determines a $u$--$v$ route.  The capacity-one arc
$x^-x^+$ permits at most one route to use an internal vertex $x$.  Conversely,
a family of internally disjoint paths gives an integral flow by sending one
unit along each corresponding route.  All finite minimum cuts can therefore
be chosen among the split arcs, and their capacities count deleted vertices.
\end{proof}

\begin{noteremark}
Taking the minimum over all relevant pairs $u,v$ recovers the vertex
connectivity $\kappa(G)$.  Adjacent pairs require the usual harmless
modification that forbids the direct edge from bypassing an internal
separator.
\end{noteremark}
\begin{noteapplication}[baseball elimination]
Suppose teams $1,\dots,n$ have current wins $w_i$, and let $r_{ij}$ be the
number of remaining games between teams $i$ and $j$.  To test whether team
$n$ can still finish tied for first, put
\[
  W:=w_n+\sum_{i<n}r_{in},
\]
its largest possible final number of wins.  If some $w_i>W$, team $n$ is
already eliminated.  Otherwise build a network with
\begin{itemize}
  \item a game vertex $x_{ij}$ for every $1\le i<j<n$;
  \item a team vertex $y_i$ for every $i<n$;
  \item arcs $s x_{ij}$ of capacity $r_{ij}$;
  \item arcs $x_{ij}y_i$ and $x_{ij}y_j$ of infinite capacity;
  \item arcs $y_i t$ of capacity $W-w_i$.
\end{itemize}
Team $n$ is not eliminated if and only if a maximum flow saturates every arc
leaving $s$.
\end{noteapplication}
\begin{proof}
A unit sent from $x_{ij}$ to $y_i$ assigns one remaining game between $i$ and
$j$ as a win for team $i$.  Saturating all source arcs assigns every remaining
game not involving team $n$, while the capacity $W-w_i$ prevents team $i$
from exceeding team $n$'s best possible total.  The correspondence is
reversible.
\end{proof}
\begin{notedefinition}[circulation with demands]
Let $D=(V,A)$ have capacities $c\colon A\to\R_{\ge0}$ and demands
$d\colon V\to\R$.  Interpret $d(v)>0$ as net demand and $d(v)<0$ as net
supply.  A \emph{circulation satisfying $d$} is a function
$f\colon A\to\R_{\ge0}$ such that
\[
  f(e)\le c(e)\quad(e\in A),
  \qquad
  f_{\inflow}(v)-f_{\out}(v)=d(v)\quad(v\in V).
\]
Necessarily $\sum_{v\in V}d(v)=0$.
\end{notedefinition}

\begin{noteapplication}[reducing circulation to maximum flow]
Assume $\sum_v d(v)=0$.  Add a super-source $s$ and a super-sink $t$.
For each supply vertex $v$ with $d(v)<0$, add $sv$ of capacity $-d(v)$.
For each demand vertex $v$ with $d(v)>0$, add $vt$ of capacity $d(v)$.
Retain all original arcs and capacities.  A feasible circulation exists if
and only if a maximum $s$--$t$ flow saturates every arc leaving $s$ (and hence
every arc entering $t$).
\end{noteapplication}
\begin{proof}
Given a feasible circulation, append the prescribed amounts on the new arcs.
The demand equation becomes ordinary flow conservation at every original
vertex, so the resulting flow saturates all source arcs.  Conversely, delete
the new arcs from any flow that saturates them.  Flow conservation then gives
$f_{\inflow}(v)-f_{\out}(v)=d(v)$ at each original vertex.
\end{proof}

\chapter{Matchings}

\section{Basic definitions and augmenting paths}

\begin{notedefinition}
Let $G=(V,E)$ be a graph.
\begin{itemize}
  \item A \emph{matching} is a set $M\subseteq E$ whose edges have pairwise
        disjoint endpoint sets.
  \item A matching is \emph{maximal} if no edge can be added to it, and
        \emph{maximum} if it has largest possible cardinality.
  \item A vertex incident with an edge of $M$ is \emph{saturated} by $M$;
        otherwise it is \emph{unsaturated} or \emph{exposed}.
  \item A \emph{perfect matching} saturates every vertex.
\end{itemize}
The maximum size of a matching is denoted $\nu(G)$.
\end{notedefinition}

\begin{notedefinition}
Let $M$ be a matching.  A path is \emph{$M$-alternating} if its edges
alternate between $M$ and $E\setminus M$.  It is \emph{$M$-augmenting} if
both endpoints are unsaturated by $M$.  Every augmenting path has odd length.
\end{notedefinition}

\begin{notedefinition}
For sets $A$ and $B$, their \emph{symmetric difference} is
\[
  A\triangle B=(A\setminus B)\cup(B\setminus A).
\]
\end{notedefinition}
\begin{notetheorem}
If $G$ is connected and $\Delta(G)\le2$, then $G$ is a path or a cycle.
\end{notetheorem}
\begin{proof}
If every vertex has degree $2$, start at any vertex and follow unused edges.
Because the graph is finite and no vertex permits branching, the first repeat
closes a cycle.  Connectedness then forces that cycle to contain every vertex.

Otherwise $G$ has a vertex $v$ of degree at most $1$.  Starting from $v$ and
extending a path as far as possible produces no repeated vertex, since a repeat
would create a cycle and leave no available edge through which the path could
have entered or left it.  If a vertex lay outside the resulting maximal path,
connectedness would give an edge from the path to the outside and hence a
further extension.  Thus the path contains all vertices.
\end{proof}

\begin{notetheorem}[Berge's Theorem]
A matching $M$ is maximum if and only if there is no $M$-augmenting path.
\end{notetheorem}
\begin{proof}
If $P$ is $M$-augmenting, then
\[
  M':=M\triangle E(P)
\]
is a matching and $|M'|=|M|+1$, so $M$ is not maximum.

Conversely, suppose a matching $N$ satisfies $|N|>|M|$.  In the spanning
subgraph with edge set $M\triangle N$, every vertex has degree at most $2$.
Its components are therefore paths and even cycles whose edges alternate
between $M$ and $N$.  Since $N$ has more edges overall, some path component
contains one more edge of $N$ than of $M$.  That path starts and ends with
$N$-edges, so both endpoints are unsaturated by $M$; it is an
$M$-augmenting path.
\end{proof}

\begin{noteremark}
Berge's theorem turns maximum matching into a search problem: repeatedly find
an augmenting path and take the symmetric difference.  In bipartite graphs,
this search can also be implemented as a network-flow computation.
\end{noteremark}
\section{Hall's theorem and its consequences}

\begin{notedefinition}
A graph is \emph{$k$-regular} if every vertex has degree $k$.  It is
\emph{regular} if it is $k$-regular for some $k\ge0$.
\end{notedefinition}

\begin{noteproposition}
If $k\ge1$, every $k$-regular bipartite graph $G=(X,Y;E)$ satisfies
$|X|=|Y|$.
\end{noteproposition}
\begin{proof}
Counting edges from both sides gives
\[
  k|X|=\sum_{x\in X}\deg(x)=|E|
  =\sum_{y\in Y}\deg(y)=k|Y|.
\]
\end{proof}

\begin{notetheorem}
Every $k$-regular bipartite graph with $k\ge1$ has a perfect matching.
\end{notetheorem}
\begin{proof}
This follows from Hall's theorem below.  For $S\subseteq X$, count the edges
between $S$ and $N(S)$:
\[
  k|S|
  =e(S,N(S))
  \le k|N(S)|.
\]
Thus $|S|\le|N(S)|$, so there is a matching saturating $X$.  Since
$|X|=|Y|$, it is perfect.
\end{proof}
\begin{notedefinition}
For $S\subseteq V(G)$, the \emph{neighbourhood} of $S$ is
\[
  N(S):=\{v\in V(G):uv\in E(G)\text{ for some }u\in S\}.
\]
A matching is \emph{$S$-saturating} if it saturates every vertex of $S$.
\end{notedefinition}

\begin{notetheorem}[Hall's Marriage Theorem]
Let $G=(X,Y;E)$ be bipartite.  There is a matching saturating $X$ if and only
if
\[
  |S|\le|N(S)|\qquad\text{for every }S\subseteq X.
\]
\end{notetheorem}
\begin{proof}
Necessity is immediate: distinct vertices of $S$ must be matched to distinct
vertices of $N(S)$.

For sufficiency, induct on $|X|$.  The result is trivial when $|X|=1$.
Assume $|X|>1$.

First suppose every nonempty proper subset $S\subsetneq X$ satisfies the
strict inequality $|N(S)|\ge|S|+1$.  Choose an edge $xy$.  In
$G-\{x,y\}$, for every $S\subseteq X\setminus\{x\}$,
\[
  |N_{G-\{x,y\}}(S)|\ge |N_G(S)|-1\ge |S|.
\]
Induction gives a matching saturating $X\setminus\{x\}$, and adding $xy$
finishes.

Otherwise choose a nonempty proper $S\subsetneq X$ with $|N(S)|=|S|$.
Hall's condition holds in the induced bipartite graph $G[S,N(S)]$, so by
induction it has a matching saturating $S$.  It also holds in
$G[X\setminus S,Y\setminus N(S)]$: for $R\subseteq X\setminus S$,
\[
  |N(R)\setminus N(S)|
  =|N(R\cup S)|-|N(S)|
  \ge |R\cup S|-|S|=|R|.
\]
Induction gives a matching saturating $X\setminus S$.  The two matchings are
vertex-disjoint, and their union saturates $X$.
\end{proof}
\begin{noteapplication}
Theorem~6.4 is an immediate application of Hall's theorem: every
$k$-regular bipartite graph with $k\ge1$ has a perfect matching.
\end{noteapplication}

\begin{notedefinition}
Let $S_1,\dots,S_n$ be finite sets.  A \emph{system of distinct
representatives} is a choice of elements $x_i\in S_i$ such that
$x_1,\dots,x_n$ are pairwise distinct.
\end{notedefinition}

\begin{notetheorem}[Hall's theorem for set systems]
The family $(S_1,\dots,S_n)$ has a system of distinct representatives if and
only if
\[
  |I|\le\left|\bigcup_{i\in I}S_i\right|
  \qquad\text{for every }I\subseteq\{1,\dots,n\}.
\]
\end{notetheorem}
\begin{proof}
Build a bipartite graph whose left vertices are the indices $1,\dots,n$ and
whose right vertices are the elements appearing in the sets, joining $i$ to
$x$ exactly when $x\in S_i$.  A matching saturating the left side is exactly
a system of distinct representatives, and Hall's condition becomes the
displayed inequality.
\end{proof}
\begin{notedefinition}
An $m\times n$ \emph{Latin rectangle} is an $m\times n$ matrix with entries
from an $n$-element symbol set such that no symbol is repeated in any row or
column.  A Latin square is an $n\times n$ Latin rectangle.
\end{notedefinition}

\begin{notetheorem}
Every $m\times n$ Latin rectangle can be extended to an $n\times n$ Latin
square by adding $n-m$ rows.
\end{notetheorem}
\begin{proof}
To add one row, form a bipartite graph whose left side consists of the $n$
columns and whose right side consists of the $n$ symbols.  Join column $j$ to
symbol $a$ exactly when $a$ has not yet appeared in column $j$.  Every column
is missing $n-m$ symbols.  Every symbol has appeared once in each of the $m$
rows, hence in $m$ distinct columns, so it too has degree $n-m$.  The graph is
regular bipartite and therefore has a perfect matching.  Put in each column
the symbol assigned by this matching.  This creates a new valid row.  Repeat
until $n$ rows have been obtained.
\end{proof}
\section{Doubly stochastic matrices}

\begin{notedefinition}
A square matrix is \emph{doubly stochastic} if all entries are nonnegative and
every row sum and every column sum is $1$.  A \emph{permutation matrix} is a
$0$--$1$ matrix with exactly one $1$ in every row and every column.
\end{notedefinition}

\begin{notedefinition}
A \emph{convex combination} of objects $v_1,\dots,v_k$ in a real vector space
is an expression
\[
  \lambda_1v_1+\cdots+\lambda_kv_k,
  \qquad
  \lambda_i\ge0,
  \quad
  \sum_{i=1}^k\lambda_i=1.
\]
\end{notedefinition}

\begin{notetheorem}
Every convex combination of permutation matrices is doubly stochastic.
\end{notetheorem}
\begin{proof}
Nonnegativity is preserved, and each row or column sum is the corresponding
convex combination of $1$'s, hence equals $1$.
\end{proof}

\begin{notetheorem}[Birkhoff--von Neumann Theorem]
Every doubly stochastic matrix is a convex combination of permutation
matrices.
\end{notetheorem}
\begin{proof}
Let $A=(a_{ij})$ be doubly stochastic.  Form the support graph with a row
vertex $r_i$, a column vertex $c_j$, and an edge $r_ic_j$ whenever
$a_{ij}>0$.  For any set $S$ of row vertices, all mass in those rows lies in
$N(S)$, so
\[
  |S|=\sum_{r_i\in S}\sum_j a_{ij}
  \le\sum_{c_j\in N(S)}\sum_i a_{ij}=|N(S)|.
\]
Hall's theorem gives a perfect matching, hence a permutation matrix $P$ whose
$1$-entries occur at positive entries of $A$.

Let $\mu$ be the minimum selected entry.  If $\mu=1$, then the row and column
sum conditions force $A=P$.  Otherwise
\[
  B:=\frac{A-\mu P}{1-\mu}
\]
is again doubly stochastic and has fewer positive entries than $A$.  Induction
on the number of positive entries writes $B$ as a convex combination of
permutation matrices, and
\[
  A=\mu P+(1-\mu)B
\]
gives the required representation.
\end{proof}
\section{Vertex covers and stable matchings}

\begin{notedefinition}
A \emph{vertex cover} is a set $C\subseteq V(G)$ incident with every edge.
The minimum size of a vertex cover is denoted $\tau(G)$.
\end{notedefinition}

\begin{notetheorem}[K\H{o}nig--Egerv\'ary Theorem]
For every bipartite graph $G$,
\[
  \nu(G)=\tau(G).
\]
\end{notetheorem}
\begin{proof}
Every vertex cover meets every edge of a matching, so $\nu(G)\le\tau(G)$.
For the reverse inequality, let $G=(X,Y;E)$ and let $M$ be a maximum matching.
Let $U\subseteq X$ be the vertices unsaturated by $M$, and let $Z$ be the set
of vertices reachable from $U$ by $M$-alternating paths that begin with an
edge outside $M$.  Put $Z_X=Z\cap X$ and $Z_Y=Z\cap Y$.

There is no augmenting path, so every vertex of $Z_Y$ is matched.  Its matched
partner lies in $Z_X\setminus U$, and this gives a bijection
$Z_Y\leftrightarrow Z_X\setminus U$.  Hence
\[
  |Z_Y|=|Z_X|-|U|.
\]
Now
\[
  C:=(X\setminus Z_X)\cup Z_Y
\]
is a vertex cover.  Indeed, an edge from a vertex of $Z_X$ to a vertex outside
$Z_Y$ would extend an alternating reachability path.  Finally,
\[
  |C|=|X|-|Z_X|+|Z_Y|=|X|-|U|=|M|.
\]
Thus $\tau(G)\le|C|=\nu(G)$.
\end{proof}

\begin{noteremark}
Under the standard matching-to-flow reduction, the minimum vertex cover in
K\H{o}nig's theorem corresponds to a minimum cut.
\end{noteremark}
\begin{notedefinition}
Let $G=K_{n,n}$ with bipartition $(X,Y)$.  Suppose every vertex has a strict
ranking of the opposite part.  A perfect matching $M$ is \emph{stable} if
there is no unmatched pair $xy$ such that $x$ prefers $y$ to its partner in
$M$ and $y$ prefers $x$ to its partner in $M$.
\end{notedefinition}

\begin{notetheorem}
Every such preference instance has a stable perfect matching.
\end{notetheorem}

\begin{notealgorithm}[Gale--Shapley]
Initially everyone is unmatched.  While some $x\in X$ is unmatched and has
not proposed to every vertex of $Y$, let $x$ propose to the most-preferred
vertex $y$ to whom it has not yet proposed.  The vertex $y$ keeps the better
of its current proposer and $x$, rejecting the other.  When no proposal is
possible, return the held pairs.
\end{notealgorithm}
\begin{proof}[Correctness]
No ordered pair is considered twice, so the algorithm stops after at most
$n^2$ proposals.  Once a vertex of $Y$ receives a proposal, it remains matched
and can only replace its partner by a more-preferred one.  If some $x\in X$
were unmatched at termination, it would have proposed to every $y\in Y$; all
$n$ vertices of $Y$ would then be holding proposals from the remaining
$n-1$ vertices of $X$, which is impossible.  Thus the returned matching is
perfect.

Suppose an unmatched pair $xy$ blocked the returned matching.  Since $x$
prefers $y$ to its final partner, $x$ proposed to $y$ earlier.  At that time
$y$ rejected $x$, immediately or later, for a partner it preferred.  Because
$y$ only trades upward, it prefers its final partner to $x$, a contradiction.
\end{proof}
\begin{noteremark}
The bipartite hypothesis is essential for this formulation.  Stable matching
in general graphs is a different problem and requires a different algorithm.
\end{noteremark}
\section{The Tutte--Berge formula}

\begin{notedefinition}
For $S\subseteq V(G)$, let $\odd(G-S)$ be the number of odd-order components
of $G-S$, and define
\[
  \deficiency(S):=\odd(G-S)-|S|.
\]
\end{notedefinition}

\begin{notelemma}[deficiency lemma]
Every finite graph $G$ has a set $S\subseteq V(G)$ and a maximum matching $M$
such that
\begin{itemize}
  \item every even component of $G-S$ has a perfect matching;
  \item every odd component of $G-S$ has a matching missing exactly one
        vertex;
  \item $M$ matches the vertices of $S$ to distinct odd components of
        $G-S$.
\end{itemize}
Consequently, $M$ misses exactly $\odd(G-S)-|S|$ vertices.
\end{notelemma}
\begin{proof}[Proof sketch]
Choose $S$ so that $\deficiency(S)$ is maximum and, subject to this, $|S|$ is
maximum.  A standard alternating-path argument shows that every even
component of $G-S$ is factorable, every odd component is factor-critical, and
the bipartite incidence graph between $S$ and the odd components satisfies
Hall's condition.  Match $S$ into distinct odd components, use a near-perfect
matching in each selected odd component that leaves the attachment vertex
free, and use perfect or near-perfect matchings in all remaining components.
If any of these assertions failed, an alternating-path exchange would either
increase the deficiency or enlarge $S$ without decreasing it, contradicting
the choice of $S$.
\end{proof}

\begin{notetheorem}[Tutte--Berge Formula]
If $G$ has $n\ge1$ vertices, then
\[
  \nu(G)
  =\frac12\min_{S\subseteq V(G)}
    \bigl(n+|S|-\odd(G-S)\bigr)
  =\frac12\left(n-
    \max_{S\subseteq V(G)}\deficiency(S)\right).
\]
\end{notetheorem}
\begin{proof}
Let $M$ be any matching and $S\subseteq V(G)$.  At most $|S|$ odd components
of $G-S$ can send an $M$-edge to $S$.  Every remaining odd component contains
an exposed vertex.  Hence
\[
  n-2|M|\ge \odd(G-S)-|S|,
\]
or equivalently
\[
  |M|\le\frac12\bigl(n+|S|-\odd(G-S)\bigr).
\]
This proves the upper bound for every $S$.  The deficiency lemma supplies a
set $S$ and a matching that misses exactly $\odd(G-S)-|S|$ vertices, so
its size is
\[
  \frac12\bigl(n+|S|-\odd(G-S)\bigr).
\]
The upper bound is therefore attained.
\end{proof}

\begin{notecorollary}[Tutte's 1-Factor Theorem]
A graph $G$ has a perfect matching if and only if
\[
  \odd(G-S)\le|S|
  \qquad\text{for every }S\subseteq V(G).
\]
\end{notecorollary}
\begin{proof}
By the Tutte--Berge formula, $\nu(G)=|V(G)|/2$ exactly when
$\deficiency(S)\le0$ for every $S\subseteq V(G)$, which is precisely the
displayed condition.
\end{proof}
\section{Exercises and further consequences}

\begin{noteexercise}
Let $G$ be bipartite.  Then $G$ has a perfect matching if and only if
\[
  \alpha(G)=\frac{|V(G)|}{2},
\]
where $\alpha(G)$ is the maximum size of an independent set.
\end{noteexercise}
\begin{proof}
If $M$ is perfect, an independent set contains at most one endpoint of each
edge of $M$, so $\alpha(G)\le|V|/2$.  Either bipartition class is independent,
so equality holds.

Conversely, the complement of a vertex cover is independent and conversely,
so $\alpha(G)+\tau(G)=|V|$.  If $\alpha(G)=|V|/2$, then
$\tau(G)=|V|/2$.  K\H{o}nig's theorem gives
$\nu(G)=\tau(G)=|V|/2$, hence a perfect matching.
\end{proof}

\begin{noteexercise}
Every bipartite graph $G$ with at least one edge has a matching of size at
least
\[
  \frac{|E(G)|}{\Delta(G)}.
\]
\end{noteexercise}
\begin{proof}
Let $C$ be a minimum vertex cover.  Since every vertex of $C$ is incident with
at most $\Delta(G)$ edges,
\[
  |E(G)|\le\Delta(G)|C|.
\]
By K\H{o}nig's theorem, $|C|=\nu(G)$.
\end{proof}
\begin{noteexercise}
Let $T$ be a tree.
\begin{enumerate}
  \item If $T$ has a perfect matching, then it is unique.
  \item The tree $T$ has a perfect matching if and only if
        \[
          \odd(T-v)=1\qquad\text{for every }v\in V(T).
        \]
\end{enumerate}
\end{noteexercise}
\begin{proof}
For (1), if $M$ and $M'$ were distinct perfect matchings, every component of
$M\triangle M'$ would be an even cycle, impossible in a tree.

For (2), suppose $M$ is perfect and remove $v$.  Exactly one component of
$T-v$ contains the vertex matched to $v$; that component has odd order, while
every other component is perfectly matched internally and has even order.
Thus $\odd(T-v)=1$.

Conversely, assume the displayed condition and induct on $|V(T)|$.  The case
$|V(T)|=2$ is immediate.  Choose a leaf $x$ and let $y$ be its neighbour.
The singleton $\{x\}$ is an odd component of $T-y$; because
$\odd(T-y)=1$, every other component $C_1,\dots,C_k$ of $T-y$ has even
order.

For $u\in V(C_i)$, all vertices outside $C_i$ lie in one connected part of
$T-u$ attached through $y$, and that outside part has even order.  Joining it
to the relevant component of $C_i-u$ therefore does not change the parity of
that component.  Consequently
\[
  \odd(C_i-u)=\odd(T-u)=1
  \qquad (u\in V(C_i)).
\]
Each $C_i$ is smaller than $T$, so induction gives a perfect matching in every
$C_i$.  Together with the edge $xy$, these matchings form a perfect matching
of $T$.
\end{proof}
\begin{noteexercise}[Petersen's Theorem]
Every bridgeless $3$-regular graph has a perfect matching.
\end{noteexercise}
\begin{proof}
By Tutte's theorem it suffices to prove
$\odd(G-S)\le|S|$ for every $S\subseteq V(G)$.  Let $C$ be an odd component
of $G-S$.  Degree counting gives
\[
  3|V(C)|=2|E(C)|+|[V(C),S]|,
\]
so the cut $[V(C),S]$ has odd size.  It cannot have size $1$, since that edge
would be a bridge.  Hence every odd component sends at least three edges to
$S$.  Therefore
\[
  3\odd(G-S)
  \le |[S,V(G)\setminus S]|
  \le 3|S|,
\]
and Tutte's condition follows.
\end{proof}

\chapter{Planar Graphs}

\section{Embeddings and Euler's formula}

\begin{notedefinition}
An \emph{embedding} of a graph $G$ in the plane assigns distinct points to
vertices and simple arcs to edges so that an arc has precisely its two
endpoints at the corresponding vertices and two edge-arcs meet only at a
common endpoint.  A graph is \emph{planar} if it has such an embedding, and a
particular embedded copy is a \emph{plane graph}.
\end{notedefinition}

\begin{noteremark}[terminology]
The connected regions of the complement of a plane graph are its
\emph{faces}.  Exactly one face is unbounded.  The boundary walk of a face may
repeat an edge when that edge is a bridge.  Counting boundary occurrences,
every edge is incident with two face-sides, and therefore
\[
  \sum_{F}\deg(F)=2|E(G)|.
\]
\end{noteremark}
\begin{notetheorem}[Euler's Formula]
Let $G$ be a connected plane graph with $n$ vertices, $m$ edges, and $f$
faces.  Then
\[
  n-m+f=2.
\]
\end{notetheorem}
\begin{proof}
Induct on $m$.  If $G$ is a tree, then $m=n-1$ and there is one face, so the
formula holds.  Otherwise $G$ has an edge $e$ on a cycle.  Such an edge is not
a bridge, so $G-e$ remains connected.  Deleting $e$ merges the two faces
incident with it; thus $n$ is unchanged while $m$ and $f$ both decrease by
$1$.  The value of $n-m+f$ is therefore unchanged, and induction completes
the proof.
\end{proof}

\begin{notecorollary}
If a plane graph has $n$ vertices, $m$ edges, $f$ faces, and $k$ connected
components, then
\[
  n-m+f=k+1.
\]
\end{notecorollary}
\begin{proof}
Apply Euler's formula to each component, placing all components in the same
unbounded face.  Equivalently, add $k-1$ noncrossing edges joining the
components and apply Theorem~7.1.
\end{proof}
\begin{notecorollary}
Let $G$ be a simple connected planar graph with $n\ge3$ vertices and $m$
edges.
\begin{enumerate}
  \item $m\le3n-6$.
  \item If $G$ contains no triangle, then $m\le2n-4$.
\end{enumerate}
\end{notecorollary}
\begin{proof}
In any plane embedding, every face has boundary length at least $3$, so
\[
  2m=\sum_F\deg(F)\ge3f.
\]
Using $f=2-n+m$ gives $m\le3n-6$.

If $G$ is triangle-free, every face boundary has length at least $4$; hence
$2m\ge4f$, and Euler's formula gives $m\le2n-4$.
\end{proof}

\begin{notecorollary}
The graphs $K_5$ and $K_{3,3}$ are not planar.
\end{notecorollary}
\begin{proof}
For $K_5$, $n=5$ and $m=10>3n-6=9$.  The graph $K_{3,3}$ is triangle-free
with $n=6$ and $m=9>2n-4=8$.
\end{proof}

\begin{notecorollary}
Every simple planar graph has a vertex of degree at most $5$.
\end{notecorollary}
\begin{proof}
For $n\le2$ the statement is immediate.  For $n\ge3$, if every vertex had
degree at least $6$, then
\[
  2m=\sum_{v\in V(G)}\deg(v)\ge6n,
\]
whereas $m\le3n-6$ gives $2m\le6n-12$, a contradiction.
\end{proof}
\section{Maximal planar graphs and triangulations}

\begin{notedefinition}
A simple planar graph is \emph{maximal planar} if adding an edge between any
two nonadjacent vertices makes the graph nonplanar.
\end{notedefinition}

\begin{notedefinition}
A plane embedding is a \emph{triangulation} if the boundary of every face is
a $3$-cycle.
\end{notedefinition}

\begin{notetheorem}
Let $G$ be a simple planar graph with $n\ge3$ vertices and $m$ edges.  The
following are equivalent.
\begin{enumerate}
  \item $G$ is maximal planar.
  \item Every plane embedding of $G$ is a triangulation.
  \item $m=3n-6$.
\end{enumerate}
\end{notetheorem}
\begin{proof}
Assume $G$ is maximal planar.  It is connected; otherwise an edge could be
added between two components.  It has no bridge and no cut-vertex, since in
either case vertices occurring consecutively on a nontriangular face could be
joined without a crossing.  Hence every face boundary is a cycle.  If some
face had length at least $4$, two nonconsecutive boundary vertices could be
joined through the interior of that face, contradicting maximality.  Thus
every embedding is a triangulation.

If every face is a triangle, then $2m=3f$.  Together with Euler's formula this
gives $m=3n-6$.

Finally, if $m=3n-6$, adding any missing edge would produce a simple graph on
$n$ vertices with $3n-5$ edges, contradicting Corollary~7.3.  Hence $G$ is
maximal planar.
\end{proof}
\section{Subdivisions and minors}

\begin{notedefinition}
To \emph{subdivide} an edge $uv$ means to replace it by a path $uwv$ through
a new vertex $w$.  A graph obtained by a sequence of edge subdivisions is a
\emph{subdivision} of the original graph.
\end{notedefinition}

\begin{noteremark}
Planarity is preserved in both directions under subdivision.  Every subgraph
of a planar graph is planar.  Consequently, a planar graph cannot contain a
subdivision of $K_5$ or $K_{3,3}$.
\end{noteremark}

\begin{notetheorem}[Kuratowski's Theorem]
A graph is planar if and only if it contains no subgraph that is a subdivision
of $K_5$ or $K_{3,3}$.
\end{notetheorem}
\begin{proof}[Comment on the proof]
The easy direction follows from the preceding remark and the nonplanarity of
$K_5$ and $K_{3,3}$.  The converse is the substantive content of
Kuratowski's theorem; it does not follow merely from closure of planarity
under subgraphs and subdivisions.  The source's one-line argument established
only the easy direction, so the difficult converse is recorded here as the
named theorem rather than given an incomplete proof.
\end{proof}
\begin{notedefinition}
To \emph{contract} an edge $xy$ means to identify $x$ and $y$, retaining all
incident edges and then deleting any loops; parallel edges may be kept or
merged when discussing simple minors.  A graph $H$ is a \emph{minor} of $G$
if $H$ can be obtained from $G$ by vertex deletions, edge deletions, and edge
contractions.
\end{notedefinition}

\begin{notetheorem}[Wagner's Theorem]
A graph is planar if and only if it has neither $K_5$ nor $K_{3,3}$ as a
minor.
\end{notetheorem}
\begin{proof}
Deleting or contracting an edge in a planar embedding preserves planarity,
so a planar graph has no nonplanar minor.  In particular, it has neither
$K_5$ nor $K_{3,3}$ as a minor.

Conversely, if $G$ is nonplanar, Kuratowski's theorem gives a subgraph $H$
that is a subdivision of $K_5$ or $K_{3,3}$.  Delete all material outside
$H$, then contract the internally vertex-disjoint paths that replaced the
original edges.  This produces $K_5$ or $K_{3,3}$ as a minor of $G$.
\end{proof}
\section{Questions from the source}

\begin{noteexercise}
If $G$ contains a subdivision of a graph $H$, then $G$ contains $H$ as a
minor.
\end{noteexercise}
\begin{proof}
Delete all edges and vertices outside the subdivision and contract each
subdivided path to a single edge.
\end{proof}

\begin{noteexercise}
The converse of Exercise~7.9 is false: a graph may contain $K_5$ as a minor
without containing a subdivision of $K_5$.
\end{noteexercise}
\begin{proof}
Start with $K_5$ on vertices $z,a,b,c,d$.  Split $z$ into adjacent vertices
$x$ and $y$, join $x$ to $a,b$, join $y$ to $c,d$, and retain all six edges
among $a,b,c,d$.  Contracting $xy$ recovers $K_5$, so $K_5$ is a minor.
However, only $a,b,c,d$ have degree at least $4$, while a subdivision of
$K_5$ requires five branch vertices of degree at least $4$.  Hence no such
subdivision exists.
\end{proof}

\begin{noteexercise}
The Petersen graph is nonplanar.
\end{noteexercise}
\begin{proof}
The Petersen graph has $n=10$, $m=15$, is bridgeless, and has girth $5$.  In a planar embedding, every face would therefore have boundary length at least $5$, so $2m\ge5f$.  Euler's
formula gives $f=2-n+m=7$, whence $30=2m\ge35$, a contradiction.
\end{proof}

\chapter{Graph Colouring}

\section{Vertex colourings}

\begin{notedefinition}
Let $G=(V,E)$ be a graph.  A \emph{proper vertex colouring} is a map
$c\colon V\to S$ such that $c(u)\ne c(v)$ whenever $uv\in E$.  If
$|S|=k$, it is a \emph{$k$-colouring}.  The least such $k$ is the
\emph{chromatic number} $\chi(G)$.
\end{notedefinition}

\begin{notetheorem}[Four-Colour Theorem]
Every planar graph is $4$-colourable.
\end{notetheorem}

\begin{notetheorem}[Five-Colour Theorem]
Every planar graph is $5$-colourable.
\end{notetheorem}
\begin{proof}
Suppose a counterexample $G$ has as few vertices as possible, and fix a plane
embedding.  By Corollary~7.5, choose a vertex $v$ with $\deg(v)\le5$.  By
minimality, $G-v$ has a $5$-colouring.  If at most four colours occur on
$N(v)$, an unused colour extends the colouring to $v$.

It remains to consider $\deg(v)=5$ with all five colours appearing.  List the
neighbours $v_1,v_2,v_3,v_4,v_5$ in cyclic order around $v$, with
$c(v_i)=i$.  Let $H_{ij}$ be the subgraph induced by vertices of colours $i$
and $j$.

If $v_1$ and $v_3$ lie in different components of $H_{13}$, interchange
colours $1$ and $3$ in the component containing $v_1$.  Colour $1$ is then
absent from $N(v)$ and may be assigned to $v$.

Otherwise there is a $1$--$3$ Kempe path from $v_1$ to $v_3$.  Together with
$v v_1$ and $v v_3$, it forms a closed curve separating $v_2$ from $v_4$.
Hence no $2$--$4$ Kempe path joins $v_2$ to $v_4$.  Interchanging colours
$2$ and $4$ in the component containing $v_2$ frees colour $2$ for $v$.
Both cases contradict the choice of $G$.
\end{proof}
\begin{notetheorem}[greedy colouring]
For every graph $G$,
\[
  \chi(G)\le\Delta(G)+1.
\]
\end{notetheorem}
\begin{proof}
Order the vertices arbitrarily and colour them one at a time.  When a vertex
is reached, at most $\Delta(G)$ colours are forbidden by already coloured
neighbours, so one of $\Delta(G)+1$ colours remains available.
\end{proof}

\begin{notelemma}[Brooks ordering lemma]
Let $G$ be a connected graph with maximum degree $\Delta\ge3$.  If $G$ is
neither complete nor an odd cycle, then its vertices can be ordered so that
\begin{itemize}
  \item two nonadjacent vertices $a,b$ are coloured first with the same
        colour;
  \item every later vertex except the final one has a neighbour still
        uncoloured when it is coloured;
  \item the final vertex is adjacent to both $a$ and $b$.
\end{itemize}
\end{notelemma}
\begin{proof}[Proof sketch]
For a $2$-connected graph that is neither complete nor a cycle, one can choose
vertices $a,b$ at distance two such that $G-\{a,b\}$ is connected; let $v$ be
a common neighbour.  Take a spanning tree of $G-\{a,b\}$ rooted at $v$ and
order its vertices from leaves toward the root.  This gives the required
ordering.  If $G$ has a cut-vertex, apply the same construction successively
to end blocks, permuting colour names at the shared cut-vertex.  Complete end
blocks cannot have order $\Delta+1$, because their cut-vertex has an
additional neighbour outside the block.  This standard block reduction leaves
only complete graphs and odd cycles as the exceptional connected graphs.
\end{proof}

\begin{notetheorem}[Brooks's Theorem]
If a connected graph $G$ is neither complete nor an odd cycle, then
\[
  \chi(G)\le\Delta(G).
\]
\end{notetheorem}
\begin{proof}
The assertion is immediate for $\Delta(G)\le2$.  For $\Delta\ge3$, use the
ordering from the lemma.  Give the nonadjacent vertices $a,b$ the same colour.
Colour the remaining vertices greedily in the prescribed order using
$\Delta$ colours.  Every nonfinal vertex has an uncoloured neighbour, so at
most $\Delta-1$ colours are forbidden.  At the final vertex, the two
neighbours $a,b$ use the same colour, so again at most $\Delta-1$ distinct
colours appear.  Thus the colouring succeeds with $\Delta$ colours.
\end{proof}
\section{Edge colourings}

\begin{notedefinition}
A \emph{proper edge colouring} is a map $c\colon E(G)\to S$ such that edges
sharing a vertex receive different colours.  The least possible number of
colours is the \emph{chromatic index} $\chi'(G)$.
\end{notedefinition}

\begin{noteremark}
Every colour class in a proper edge colouring is a matching, and
\[
  \chi'(G)\ge\Delta(G).
\]
A graph is $2$-edge-colourable exactly when each nontrivial component is a
path or an even cycle.
\end{noteremark}

\begin{notelemma}[Vizing fan extension lemma]
Let $G$ be simple, let $uv$ be an uncoloured edge, and suppose all other edges
have been properly coloured with $\Delta(G)+1$ colours.  By rotating colours
on a maximal fan centred at $u$ and, if necessary, interchanging two colours
along a maximal alternating Kempe chain, the colouring can be extended to
$uv$.
\end{notelemma}
\begin{proof}[Proof sketch]
Build a maximal fan
$v=v_0,v_1,\dots,v_k$ at $u$, where the colour on $uv_i$ is missing at an
earlier fan vertex.  Let $\alpha$ be missing at $u$ and let $\beta$ be missing
at $v_k$.  Consider the maximal $\alpha$--$\beta$ alternating chain through
$v_k$.  If it avoids $u$, swap $\alpha$ and $\beta$ on that chain and rotate
the fan, making one colour free at both ends of the uncoloured edge.  If it
reaches $u$, the defining missing colours of the fan locate an earlier fan
vertex at which a second Kempe interchange followed by a shorter fan rotation
has the same effect.  Maximality of the fan guarantees that one of these two
operations applies.  All operations preserve properness because they merely
swap colours on a two-coloured component or move a missing colour along the
fan.
\end{proof}

\begin{notetheorem}[Vizing's Theorem]
For every simple graph $G$,
\[
  \Delta(G)\le\chi'(G)\le\Delta(G)+1.
\]
\end{notetheorem}
\begin{proof}
The lower bound is immediate at a vertex of maximum degree.  For the upper
bound, induct on $|E(G)|$.  Remove an edge $uv$ and colour $G-uv$ with
$\Delta(G)+1$ colours by induction.  The Vizing fan extension lemma extends
this colouring to $uv$.
\end{proof}

\begin{noteremark}
Simple graphs with $\chi'(G)=\Delta(G)$ are called \emph{class 1}; those with
$\chi'(G)=\Delta(G)+1$ are \emph{class 2}.  The theorem is false in this form
for multigraphs, where more than $\Delta+1$ colours may be needed.
\end{noteremark}
\begin{notetheorem}[K\H{o}nig's Line-Colouring Theorem]
If $G$ is bipartite, then
\[
  \chi'(G)=\Delta(G).
\]
\end{notetheorem}
\begin{proof}
Let $\Delta=\Delta(G)$.  Add dummy vertices so that the two bipartition
classes have the same size, and then add edges between deficient vertices
until every vertex has degree $\Delta$.  Parallel edges may be used in this
auxiliary bipartite multigraph $R$; the total degree deficit is the same on
both sides, so the process terminates.

Every regular bipartite multigraph has a perfect matching by Hall's theorem.
Remove one perfect matching from $R$; the remaining graph is
$(\Delta-1)$-regular.  Repeating decomposes $E(R)$ into $\Delta$ perfect
matchings.  Give each matching its own colour and restrict the colouring to
$G$.  Thus $\chi'(G)\le\Delta$, while the reverse inequality always holds.
\end{proof}
\section{Exercises and extremal examples}

\begin{noteexercise}
Suppose every two odd cycles of a graph $G$ have a common vertex.  Then
$G$ is $5$-colourable.
\end{noteexercise}
\begin{proof}
If $G$ has no odd cycle, it is bipartite and hence $2$-colourable.  Otherwise
fix an odd cycle $C$.  Every odd cycle meets $C$, so $G-V(C)$ has no odd cycle
and is bipartite.  Colour $G-V(C)$ with two colours and colour $C$ with three
new colours.  Since the two colour sets are disjoint, this is a proper
$5$-colouring of $G$.
\end{proof}
\begin{noteexercise}
For the complete graph $K_n$,
\[
  \chi'(K_n)=
  \begin{cases}
    n-1,& n\text{ even},\\
    n,& n\text{ odd}.
  \end{cases}
\]
\end{noteexercise}
\begin{proof}
The lower bound $\chi'(K_n)\ge n-1$ is the degree bound.  If $n$ is odd, each
matching contains at most $(n-1)/2$ edges, while
$|E(K_n)|=n(n-1)/2$; therefore at least $n$ colours are needed.

For odd $n$, label the vertices by $\mathbb Z_n$ and colour the edge $ij$ by
$i+j\pmod n$.  At a fixed vertex $i$, the values $i+j$ are distinct as
$j$ varies, so this is a proper $n$-edge-colouring.

For even $n=2m$, label the vertices by
$\{\infty\}\cup\mathbb Z_{2m-1}$.  For each
$c\in\mathbb Z_{2m-1}$, let
\[
  M_c
  :=\{\infty c\}
    \cup\bigl\{(c+i)(c-i):1\le i\le m-1\bigr\},
\]
with arithmetic modulo $2m-1$.  Each $M_c$ is a perfect matching, and the
$2m-1$ matchings partition $E(K_{2m})$.  Assigning one colour to each $M_c$
gives an $(n-1)$-edge-colouring.
\end{proof}
\begin{noteexercise}
Let $G$ have $n$ vertices and let $\overline G$ be its complement.  Then
\[
  \chi(G)\,\chi(\overline G)\ge n.
\]
\end{noteexercise}
\begin{proof}
Choose proper colourings
$p\colon V(G)\to[\chi(G)]$ and
$q\colon V(G)\to[\chi(\overline G)]$.  Assign each vertex the ordered pair
$(p(v),q(v))$.  Distinct vertices cannot receive the same pair: if they did,
they would be nonadjacent in $G$ because their $p$-colours agree and
nonadjacent in $\overline G$ because their $q$-colours agree, which is
impossible.  The assignment is injective, so the number of available pairs is
at least $n$.
\end{proof}


\end{document}
